\documentclass[a4paper,11pt]{article}
\pdfoutput=1
\usepackage[T1]{fontenc}
\usepackage[utf8]{inputenc}
\usepackage{lmodern}
\usepackage[a4paper,margin=1in]{geometry}
\usepackage{amsmath,amssymb,amsthm,mathtools}
\usepackage{booktabs,array,microtype}
\usepackage[numbers,sort&compress]{natbib}
\PassOptionsToPackage{hyphens}{url}
\usepackage[hidelinks]{hyperref}

\hypersetup{
  pdftitle={Further results on the square substitution: complete graphs, maps from spheres, and small graphs},
  pdfauthor={Julius A. Zeiss}
}

\newtheorem{theorem}{Theorem}[section]
\newtheorem{lemma}[theorem]{Lemma}
\newtheorem{proposition}[theorem]{Proposition}
\newtheorem{corollary}[theorem]{Corollary}
\theoremstyle{definition}
\newtheorem{definition}[theorem]{Definition}
\newtheorem{remark}[theorem]{Remark}

\newcommand{\C}{\mathbb C}
\newcommand{\R}{\mathbb R}
\newcommand{\Ftwo}{\mathbb F_2}
\newcommand{\RP}{\mathbb{RP}}
\newcommand{\cq}{\chi_{\mathrm q}}
\newcommand{\B}{\operatorname{Box}}
\newcommand{\qto}{\xrightarrow{\mathrm q}}
\DeclareMathOperator{\dist}{dist}
\DeclareMathOperator{\rank}{rank}
\DeclareMathOperator{\ran}{ran}
\DeclareMathOperator{\supp}{supp}
\DeclareMathOperator{\tr}{tr}
\numberwithin{equation}{section}

\allowdisplaybreaks[2]
\setlength{\emergencystretch}{2em}

\title{Further results on the square substitution:\\ complete graphs, maps from spheres, and small graphs}
\author{Julius A. Zeiss\thanks{Institute for Quantum Information, RWTH Aachen University, Otto-Blumenthal-Stra\ss e 20, 52074 Aachen, Germany. \texttt{jzeiss@physik.rwth-aachen.de}}}
\date{27 November 2026}

\begin{document}
\maketitle

\begin{abstract}
The article \emph{Explicit graphs with quantum chromatic number three and large chromatic number} builds its graphs by iterating one operation, the square substitution, on the triangle. This note collects results of the same project that the article does not contain. For every $n\ge3$, one square substitution over the complete graph $K_n$ has quantum chromatic number $n$ and chromatic number exactly $2n-1$. If the box complex of a graph receives an equivariant map from a sphere of dimension $2h+1$, the first substitution raises the length of the tower that certifies the chromatic number from $h$ to $2h+1$; for complete graphs this gives $\chi\ge2n-3+2j$ after $j$ substitutions for odd $n$ and $\chi\ge2n-4+2j$ for even $n$, and for even $n$ we show that at the first level the certificate cannot reach the exact value when the wires have at least two interior layers. The box complex of a graph with a quantum $3$-coloring receives no equivariant map from the two-sphere. We also describe further small graphs with quantum chromatic number three: graphs with chromatic number four and a quantum $3$-coloring whose measurements commute up to a prescribed distance, on $22+69\lceil(R-1)/2\rceil$ vertices; a graph on $165$ measurements with chromatic number five and only $678$ edges; graphs with chromatic number five on $234+54(R-2)$ vertices with a quantum $3$-coloring whose compatibility radius is exactly $R\ge2$; and a family of graphs of measurements with chromatic number at least $2j+3$ whose first level has $165$ vertices. A last section lists facts that are known only by computation, with the programs that verify them.
\end{abstract}

\section{Introduction}\label{sec:intro}

The article~\cite{Z} constructs explicit graphs with quantum chromatic number three and arbitrarily large chromatic number by iterating one graph operation, the \emph{square substitution} $S_\lambda$, on the triangle. This note collects results of the same project that are not contained in~\cite{Z}. It is a supplement to that article and is not meant to be read without it: it uses the definitions of~\cite{Z} and several of its statements, all of which Section~\ref{sec:background} restates. Numbers of statements of~\cite{Z} refer to the version of 26 November 2026.

\paragraph{Complete graphs and maps from spheres (Section~\ref{sec:complete}).}
For the triangle, one substitution has chromatic number five~\cite[Theorem~2.4]{Z}, and for every graph $H$ with $\chi(H)\ge3$ one has $\chi(S_\lambda(H))\le2\chi(H)-1$~\cite[Proposition~A.1]{Z}. Theorem~\ref{thm:complete} shows that this bound is attained by every complete graph: $\chi(S_\lambda(K_n))=2n-1$ and $\cq(S_\lambda(K_n))=n$ for all $n\ge3$. The lower bound uses a lemma of Brakensiek and Guruswami on colorings of powers of $K_n$ with $2n-2$ colors~\cite{BG}.

The lower bounds of~\cite{Z} for iterated substitutions rest on towers of chains: a tower of length $h$ over a graph certifies $2h+3$ colors, and each substitution lengthens a tower by one. A sphere $S^{2h+1}$ with its standard circle action carries a tower of length $h$. Theorem~\ref{thm:spheres} shows that after the first substitution over this sphere the tower has length $2h+1$, not $h+1$. Consequently, if the box complex of $H$ receives an equivariant map from $S^{2h+1}$, then $\chi(S^j_\lambda(H))\ge4h+2j+3$, with one more color if the sphere has dimension $2h+2$. For complete graphs this gives $\chi(S^j_\lambda(K_n))\ge2n-3+2j$ for odd $n$ and $\ge2n-4+2j$ for even $n$ (Corollary~\ref{cor:kn}), on sparse induced subgraphs as in~\cite{Z}. For odd $n$ and $j=1$ this is the exact value $2n-1$. For even $n$ it is one less, and Proposition~\ref{prop:even} shows that this is not a defect of the proof: on $S_\lambda(K_n)$ with $\lambda\ge2$ the powers of the Stiefel--Whitney class of the box complex certify exactly $2n-2$ colors.

\paragraph{Maps from the two-sphere (Section~\ref{sec:coindex}).}
For a graph $G$ with a quantum $3$-coloring, the box complex can admit no equivariant map to a low-dimensional sphere, which is how~\cite{Z} proves that $\chi(G)$ is large. In the other direction it is restricted: Proposition~\ref{prop:coindex} shows that no sphere of dimension two maps equivariantly into the box complex of $G$. The proof combines a unitary representation built from the coloring with a theorem of Simons, Tardif and Wehlau on signatures of odd closed walks~\cite{STW}.

\paragraph{Small graphs (Section~\ref{sec:small}).}
The article exhibits a graph $F_{22}$ on $22$ vertices with $(\cq,\chi)=(3,4)$ and a graph $F_{\rm phys}$ on $165$ measurements with $(\cq,\chi)=(3,5)$. We add four constructions. Table~\ref{tab:small} lists them.
\begin{itemize}
\item Proposition~\ref{prop:stretch}: for every $R\ge2$, a graph with $\chi=4$ on $22+69\lceil(R-1)/2\rceil$ vertices with an $R$-compatible quantum $3$-coloring in local dimension four. Every quantum $3$-coloring of it has local dimension divisible by four.
\item Theorem~\ref{thm:star}: spanning subgraphs $F_\triangle\supset F_*$ of $F_{\rm phys}$ with $732$ and $678$ edges that still carry a tower of length one, hence still have $\chi=5$.
\item Theorem~\ref{thm:FR}: for every $R\ge2$, a graph $F_R$ with $\chi=5$ on $234+54(R-2)$ vertices and a quantum $3$-coloring in local dimension four whose compatibility radius is exactly $R$. For comparison, $S_{R-2}(K_3)$ has $342+81(R-2)$ vertices.
\item Theorem~\ref{thm:Hj}: a family of physical graphs $H_j$ with $H_1=F_{\rm phys}$, $\cq(H_j)=3$, $\chi(H_j)\ge2j+3$ and singly exponential order. Every quantum $3$-coloring of $H_j$ has local dimension divisible by $2\cdot4^{j-1}$ and at least $4^j$. The graph $H_2$ has $4{,}330{,}533$ vertices and $7\le\chi(H_2)\le9$.
\end{itemize}

\begin{table}[t]
\centering
\small\setlength{\tabcolsep}{5pt}
\begin{tabular}{@{}llccc@{}}
\toprule
Graph & Vertices & $\chi$ & Radius & Source\\
\midrule
$F_{22}$ & $22$ & $4$ & $1$ & \cite[Prop.~6.1]{Z}\\
stretched $F_{22}$ & $22+69\lceil(R-1)/2\rceil$ & $4$ & $R\ge2$ & Prop.~\ref{prop:stretch}\\
$F_{\rm phys}$, $F_\triangle$, $F_*$ & $165$ & $5$ & $1$ & \cite[Thm.~6.2]{Z}, Thm.~\ref{thm:star}\\
$F_R$ & $234+54(R-2)$ & $5$ & $R\ge2$ & Thm.~\ref{thm:FR}\\
$S_\lambda(K_3)$, $\lambda=\max\{R-2,0\}$ & $342+81\lambda$ & $5$ & $R\ge1$ & \cite[Thm.~2.4]{Z}\\
$H_2$ & $4{,}330{,}533$ & $7$, $8$ or $9$ & $1$ & Thm.~\ref{thm:Hj}\\
\bottomrule
\end{tabular}
\caption{Graphs with quantum chromatic number three and a quantum $3$-coloring in local dimension four ($16$ for $H_2$). The column ``Radius'' gives the values of $R$ for which an $R$-compatible quantum $3$-coloring is exhibited. All entries have analytical proofs. Section~\ref{sec:computed} adds graphs with $\chi=5$ on $207$ and $279$ vertices for $R=2$ and $R=3$ whose lower bound $\chi\ge5$ is computer-assisted.}
\label{tab:small}
\end{table}

\paragraph{Facts known only by computation (Section~\ref{sec:computed}).}
The last section lists statements that were verified by programs and have no analytical proof here, among them the criticality of $F_{22}$, the minimality of the $678$ edges of $F_*$, and the edge number of $H_2$.

\paragraph{Proofs, imports and computations.}
Every numbered statement of Sections~\ref{sec:complete}--\ref{sec:small} has an analytical proof from the definitions, from statements of~\cite{Z}, and from two published results that we import and state in full: Lemma~3.7 of~\cite{BG} (as Lemma~\ref{lem:bg}), used for Theorem~\ref{thm:complete}, and Theorem~1 of~\cite{STW} with its Lemma~4 (in Section~\ref{sec:coindex}), used for Proposition~\ref{prop:coindex}. Beyond these we use standard algebraic topology: the Serre spectral sequence, the Thom isomorphism and Poincar\'e duality with coefficients in $\Ftwo$. The program \texttt{check\_note.py} that accompanies this note replays finite facts from these sections independently; none of the proofs refers to it. The facts of Section~\ref{sec:computed} are set apart because they rest on computation alone.

\section{Definitions and statements used from the article}\label{sec:background}

This section fixes the notation and restates what is used from~\cite{Z}. All graphs are finite, simple and loopless. We write $u\sim v$ for adjacency, $N(v)$ for the set of neighbors of $v$, and $[n]=\{1,\ldots,n\}$.

\subsection{Quantum colorings}\label{sec:quantum}

A \emph{quantum $c$-coloring} of a graph $G$ in \emph{local dimension} $d$ assigns to every vertex $v$ a projective measurement $(P_{v,a})_{a\in[c]}$ on $\C^d$, that is, orthogonal projections with $\sum_aP_{v,a}=I$, such that $P_{u,a}P_{v,a}=0$ for every edge $uv$ and every color $a$ (the \emph{edge rule}). The \emph{quantum chromatic number} $\cq(G)$ is the least $c$ for which $G$ has a quantum $c$-coloring in some local dimension. A quantum coloring is \emph{$R$-compatible} if all projectors at $u$ commute with all projectors at $v$ whenever $\dist_G(u,v)\le R$, and its \emph{compatibility radius} is the largest such $R$. More generally, a \emph{quantum homomorphism} $G\qto H$ assigns to every vertex $v$ of $G$ a projective measurement $(P_{v,a})_{a\in V(H)}$ such that $P_{u,a}P_{v,b}=0$ whenever $u\sim_Gv$ and $a\not\sim_Hb$; a quantum $c$-coloring is a quantum homomorphism to $K_c$. We use four standard facts.
\begin{itemize}
\item[(Q1)] In a quantum $3$-coloring, all projectors at two adjacent vertices commute~\cite[Lemma~2]{Ji}.
\item[(Q2)] On a triangle $u,v,w$ of a graph with a quantum $3$-coloring, $P_{u,a}+P_{v,a}+P_{w,a}=I$ for every color $a$. Indeed, the three projections are pairwise orthogonal, so their sum is at most $I$, and summing over the three colors gives $3I$ on both sides.
\item[(Q3)] A graph with a quantum $2$-coloring is bipartite~\cite[Proposition~3]{CMNSW}. Hence a graph that contains a triangle and has a quantum $3$-coloring has quantum chromatic number three.
\item[(Q4)] A graph that contains a clique on $n$ vertices has $\cq\ge n$: in a quantum $c$-coloring in local dimension $d$ the projectors of one color at the vertices of the clique are pairwise orthogonal, so their ranks sum to at most $d$, and summing over the colors gives $nd\le cd$.
\end{itemize}

\subsection{The six bases}\label{sec:square}

Each entry of the following table lists the coordinates of an unnormalized vector of $\C^4$, with $+$ and $-$ for $1$ and $-1$. Row $i$ is an orthogonal basis, called \emph{basis $i$}, and $Q^i_s$ denotes the projection onto the ray of the vector in column $s$. The index $s\in\{0,1,2,3\}$ is called a \emph{position}.
\[
\begin{array}{c|cccc}
 i\backslash s&0&1&2&3\\ \hline
 1& ++++ & +-+- & ++-- & +--+\\
 2& +000 & 0+00 & 00+0 & 000+\\
 3& +++- & ++-+ & +-++ & -+++\\
 4& +0+0 & +0-0 & 0+0+ & 0+0-\\
 5& ++00 & +-00 & 00++ & 00+-\\
 6& +00+ & +00- & 0++0 & 0+-0
\end{array}
\]
For $r\in\{1,2,3\}$ put $A_r=\{0,r\}$ and $P^i_r=Q^i_0+Q^i_r$. These projections satisfy the nine identities
\begin{equation}\label{eq:planes}
\begin{array}{ccc}
 P^1_1=P^4_2,&P^1_2=P^5_2,&P^1_3=P^6_2,\\
 P^2_1=P^5_1,&P^2_2=P^4_1,&P^2_3=P^6_1,\\
 P^3_1=P^5_3,&P^3_2=P^4_3,&P^3_3=I-P^6_3 .
\end{array}
\end{equation}
The six bases are the joint eigenbases of the rows and columns of the Mermin--Peres square, and the identities express that each of bases $1,2,3$ shares one decomposition of $\C^4$ into two planes with each of bases $4,5,6$. The bases are thus the vertices of a complete bipartite graph $K_{3,3}$ with sides $\{1,2,3\}$ and $\{4,5,6\}$. For an edge $e=ik$ of this $K_{3,3}$, always written with $i\le3<k$, let $P_e$ be the projection on the left of its identity in~\eqref{eq:planes}; its range is the \emph{plane} of $e$. For $b\in\{i,k\}$ let $A_{b,e}$ be the set of the two positions of basis $b$ whose rays span the plane of $e$. Then $A_{b,e}$ is one of $A_1,A_2,A_3$, except at the \emph{complemented edge} $e=36$, where $A_{3,36}=A_3$ and $A_{6,36}=A_3^c=\{1,2\}$. The three edges at a basis $b\le5$ use $A_1,A_2,A_3$ once each, and the three edges at basis $6$ use $A_1,A_2,A_3^c$. The nine planes are pairwise distinct and no two of them are complementary, and the $24$ rays of the six bases are pairwise distinct.

For $\{r,s,t\}=\{1,2,3\}$ the partition $\{A_r,\{s\},\{t\}\}$ of the set of positions, taken at a basis $i$, is a \emph{ternary template}, and the partition $\{A_3,A_3^c\}$ at basis $6$ is the \emph{binary template}. There are eighteen ternary templates, hence nineteen templates. The sets of a template are its \emph{parts}.

\subsection{The square substitution}\label{sec:substitution}

Let $H$ be a graph and $\lambda\ge0$ an integer. The graph $S_\lambda(H)$ is defined as follows~\cite[Section~2.3]{Z}. For each basis $i$ take a copy of the set of \emph{words} $g=(g_0,g_1,g_2,g_3)\in V(H)^4$ with $g_0=g_r$ for some $r\in\{1,2,3\}$, the \emph{basis chart} at $i$; two words of one chart are adjacent when their entries are adjacent in $H$ at all four positions. Along each edge $e=ik$ of $K_{3,3}$ place a \emph{wire} with \emph{layers} $0,\ldots,\lambda+1$, each a copy of $V(H)^2$. The pair $(x,y)$ of layer $0$ is identified with the word of chart $i$ that takes $x$ on $A_{i,e}$ and $y$ on $A_{i,e}^c$, and the pair $(x,y)$ of layer $\lambda+1$ with the word of chart $k$ that takes $x$ on $A_{k,e}$ and $y$ on $A_{k,e}^c$. The vertices of the $\lambda$ \emph{interior} layers are written $(e,\ell,x,y)$. Within a wire, $(x,y)$ and $(x',y')$ in equal or consecutive layers are adjacent when $x\sim x'$ and $y\sim y'$. The $j$-fold iterate is $S^j_\lambda(H)$, and its vertices are said to be at \emph{level} $j$. We call the $|V(H)|^2$ words of a chart that are identified with a layer of a wire the \emph{end grid} of that wire in that chart.

Given a quantum homomorphism from $H$ to some graph with projectors $P_{g,a}$ on $\C^d$, put
\begin{equation}\label{eq:strategy}
 P'_{(i,g),a}=\sum_{s=0}^3Q^i_s\otimes P_{g_s,a},\qquad
 P'_{(e,\ell,x,y),a}=P_e\otimes P_{x,a}+(I-P_e)\otimes P_{y,a}
\end{equation}
on $\C^4\otimes\C^d$. By~\eqref{eq:planes} the two formulas agree on the identified vertices, and they define a quantum homomorphism from $S_\lambda(H)$ to the same graph. If the given homomorphism is $R$-compatible and $R\le\lambda+2$, the new one is $R$-compatible~\cite[Lemma~2.2]{Z}. Starting from the measurements on $\C^1$ that return the vertex itself and applying~\eqref{eq:strategy} $j$ times gives a $(\lambda+2)$-compatible quantum homomorphism $S^j_\lambda(H)\qto H$ in local dimension $4^j$. For $H=K_n$ and $j=1$ the \emph{measurement} of the word $(i,g)$ assigns the label $g_s$ to the ray at position $s$ of basis $i$, and the measurement of $(e,\ell,x,y)$ assigns $x$ to the plane of $e$ and $y$ to its complement.

A graph whose vertices are distinct projective measurements and whose edges are the pairs of measurements that are orthogonal at every outcome is called \emph{physical}. In contrast, $S_\lambda(H)$ is \emph{tagged}: distinct vertices may carry the same measurement.

\paragraph{Presentations.}
A \emph{presentation} of a vertex of $S^g_\lambda(H)$ is a rooted tree~\cite[Sections~3.1 and~7]{Z}. A \emph{leaf} carries a vertex $a$ of $H$; at level $0$ it evaluates to $a$, and at level $g\ge1$ to the constant word in basis $6$ on its value at level $g-1$ (the vertex $a$ is \emph{padded}). An \emph{internal node} at level $g\ge1$ is either a \emph{basis node}, which specifies a template, has one child for each part, and evaluates to the word that repeats the value of each child at level $g-1$ on its part, or a \emph{wire node}, which specifies an edge $e$ of $K_{3,3}$ and a layer $\ell$, has a first child that is any presentation and a second child, its \emph{background}, that is a leaf, and evaluates to the pair of their values in layer $\ell$ of the wire of $e$. For integers $j,s\ge1$ we write
\[
 G^{(s)}_{j,\lambda}(H)
\]
for the subgraph of $S^j_\lambda(H)$ induced by the vertices that admit a presentation with at most $2s-1$ internal nodes and at most $s$ wire nodes. The sparse family of~\cite{Z} is $G_{j,R}=G^{(j)}_{j,\lambda}(K_3)$ with $\lambda=\max\{R-2,0\}$, and the graph $G_{j,R}(H)$ in the proof of~\cite[Theorem~7.1]{Z} is $G^{(j)}_{j,\lambda}(H)$.

\subsection{The box complex and towers}\label{sec:box}

The \emph{box complex} $\B(H)$ is the space of pairs $(\mu,\nu)$ of probability measures on $V(H)$ such that every vertex of $\supp\mu$ is adjacent to every vertex of $\supp\nu$. The two supports are the \emph{shores} of the point. Exchanging $\mu$ and $\nu$ is a free involution, a map between spaces with involutions is \emph{equivariant} if it commutes with them, and a graph homomorphism induces an equivariant map of box complexes. For a pair $(A,B)$ of nonempty sets of vertices with every vertex of $A$ adjacent to every vertex of $B$, the points with $\supp\mu\subseteq A$ and $\supp\nu\subseteq B$ form a product of two simplices; these \emph{cells} make $\B(H)$ a regular cell complex. The complex $\B(K_p)$ is equivariantly homeomorphic to the sphere $S^{p-2}$ with the antipodal involution, by $(\mu,\nu)\mapsto\mu-\nu$.

A free involution $T$ on a compact polyhedron $Z$ defines a double cover $Z\to Z/T$ and with it a class $w_Z\in H^1(Z/T;\Ftwo)$, its first Stiefel--Whitney class, which is pulled back by equivariant maps. If $Z$ maps equivariantly to $\B(H)$ and $w_Z^k\ne0$, then
\begin{equation}\label{eq:obstruction}
 \chi(H)\ge k+2 ,
\end{equation}
because a proper $p$-coloring gives an equivariant map $\B(H)\to S^{p-2}$ and $H^k(\RP^{p-2};\Ftwo)=0$ for $k>p-2$. We write $w$ for $w_{\B(H)}$ and call the largest $k$ with $w^k\ne0$ the \emph{height} of $w$.

\paragraph{The substitution on spaces.}
For a compact polyhedron $X$, the space $S(X)$ consists of six \emph{basis blocks}, the block of basis $i$ being a copy of the set of tuples $(x_0,x_1,x_2,x_3)\in X^4$ with $x_0=x_r$ for some $r\in\{1,2,3\}$, and of nine cylinders $X^2\times[0,1]$, one for each edge $e=ik$ of $K_{3,3}$, where $(x,y,0)$ is identified with the tuple of block $i$ that has $x$ on $A_{i,e}$ and $y$ on $A_{i,e}^c$, and $(x,y,1)$ with the tuple of block $k$ that has $x$ on $A_{k,e}$ and $y$ on $A_{k,e}^c$. This is the definition of~\cite[Section~3.3]{Z}, which is stated there for spaces with a circle action; it uses no action. A continuous map $X\to X'$ induces a map $S(X)\to S(X')$, and an action of a group on $X$ induces one on $S(X)$, which is free if the given action is free. Each template gives a \emph{template map} $X^3\to S(X)$ or $X^2\to S(X)$ that repeats its arguments on the parts.

A \emph{circle space} is a compact polyhedron with a free action of the circle group $S^1\subset\C$; its \emph{half-turn} is the action of $-1$. If $X$ is a circle space, so is $S(X)$. On a circle space, chains are normalized singular chains with coefficients in $\Ftwo$, and $Bc$ denotes the chain swept out by a chain $c$ under one full rotation; one has $\partial B=B\partial$ and $B^2=0$~\cite[Appendix~B.1]{Z}. A \emph{tower of length $h$} on a circle space $Z$ is a point $x_0$, regarded as a $0$-chain, with chains $x_1,\ldots,x_h$, $x_m$ of dimension $2m$, such that
\begin{equation}\label{eq:tower}
 \partial x_m=Bx_{m-1}\qquad(1\le m\le h).
\end{equation}
A \emph{tower of length $h$ over a graph $G$} is a tower of length $h$ on a circle space $Z$ together with a map $Z\to\B(G)$ that is equivariant for the half-turn. The article proves the following.
\begin{itemize}
\item[(T1)] \cite[Lemma~B.2]{Z}. If a circle space $Z$ carries a tower of length $h$, then $w_Z^{2h+1}\ne0$. Hence a tower of length $h$ over $G$ gives $w^{2h+1}\ne0$ on $\B(G)$ and $\chi(G)\ge2h+3$~\cite[Proposition~3.2]{Z}.
\item[(T2)] \cite[Lemma~B.1(a)]{Z}. If a circle space $X$ carries a tower $x_0,\ldots,x_h$, then $S(X)$ carries a tower $y_0,\ldots,y_{h+1}$ whose point $y_0$ is the constant tuple on $x_0$ in the block of basis $6$.
\item[(T3)] \cite[Lemma~B.3]{Z}. An equivariant map $\phi:X\to\B(H)$ induces an equivariant map $S(X)\to\B(S_\lambda(H))$. Its proof uses only the involution of $X$, so the statement holds for every compact polyhedron $X$ with a free involution. The induced map has the following property. Write $\phi(x)=(\mu_x,\nu_x)$. At a point $(x_0,x_1,x_2,x_3)$ of the block of basis $i$, every word in the first shore of the image has its entry at position $s$ in $\supp\mu_{x_s}$, with equal entries at positions where the points $x_s$ are equal, and likewise for the second shore with $\nu$. At a point $(x,y,t)$ of the cylinder of $e$, every vertex in the first shore is a pair in $\supp\mu_x\times\supp\mu_y$ in one of the at most two layers nearest to $(\lambda+1)t$, and likewise for the second shore.
\item[(T4)] \cite[Theorem~3.3]{Z}. By (T2) and (T3), a tower of length $h$ over $H$ on $X$ gives a tower of length $h+1$ over $S_\lambda(H)$ on $S(X)$.
\end{itemize}

\subsection{Further statements of the article}\label{sec:imports}

\begin{itemize}
\item[(Z1)] \cite[Proposition~A.1]{Z}. Let $n\ge3$ and let $\mathcal P_n$ be the physical graph on all measurements that assign labels in $[n]$ to the four rays of one of the six bases. Then $\chi(\mathcal P_n)\le2n-1$. Consequently $\chi(S_\lambda(H))\le2\chi(H)-1$ for every graph $H$ with $\chi(H)\ge3$.
\item[(Z2)] \cite[Lemma~4.2 and Facts~3 and~4 of Section~4.1]{Z}. Let $(P_{u,a})$ be a quantum $3$-coloring of the categorical power $K_3^{\times m}$ on $\C^d$. Then all its projectors commute, and there is an orthogonal decomposition $\C^d=V_1\oplus\cdots\oplus V_m$ that reduces every projector, such that on $V_i$ the measurement at $u$ depends only on $u_i$ and satisfies the edge rule of $K_3$ as a function of $u_i$. For a set $B$ of coordinates and any two words $u,u'$ that agree on $B$ and differ at every coordinate outside $B$, the operator $\sum_aP_{u,a}P_{u',a}$ is the projection onto $\bigoplus_{i\in B}V_i$, the \emph{coordinate projection} of $B$.
\item[(Z3)] \cite[Lemma~4.3]{Z}. Let $U$ be a $3\times3$ array of self-adjoint involutions on a finite-dimensional space such that the entries of each row and of each column commute, all row products are $+I$ and all column products are $-I$. Then for two distinct entries in a common row or column, the four joint eigenspaces have equal dimension.
\item[(Z4)] \cite[Theorem~5.5(a)]{Z}. Every graph with a quantum $3$-coloring in local dimension one or two is classically $3$-colorable.
\item[(Z5)] \cite[Proposition~5.8]{Z}. If $G$ has a quantum $3$-coloring in local dimension $d=2^\nu m$ with $m$ odd, then $w^{2^{\nu+1}}=0$ on $\B(G)$. In particular, $w^2\ne0$ forces every quantum $3$-coloring of $G$ to have even local dimension.
\item[(Z6)] \cite[Proposition~6.1 and Proposition~C.1]{Z}. The graph $F_{22}$, recalled in Section~\ref{sec:stretch}, has $22$ vertices, $42$ edges, $\cq=3$ and $\chi=4$, and every quantum $3$-coloring of it has local dimension divisible by four.
\item[(Z7)] \cite[Theorem~6.2 and Appendix~E]{Z}. The graph $F_{\rm phys}$ on $165$ measurements, recalled in Section~\ref{sec:165}, has $(\cq,\chi)=(3,5)$ and carries a tower of length one.
\end{itemize}

\section{Complete graphs and maps from spheres}\label{sec:complete}

\subsection{Colorings of powers of \texorpdfstring{$K_n$}{Kn} with \texorpdfstring{$2n-2$}{2n-2} colors}\label{sec:dictator}

The categorical power $K_n^{\times L}$ has vertex set $[n]^L$, and two tuples are adjacent when they differ in every coordinate. We import the following lemma. In~\cite{BG} it is stated, with $t$ in place of $n$, for functions that satisfy the ``$[2,t,2t-2]$-colorability gadget'', which by~\cite[Proposition~2.4]{BG} are the proper colorings of $K_t^{\times L}$ with $2t-2$ colors.

\begin{lemma}[Brakensiek and Guruswami {\cite[Lemma~3.7]{BG}}]\label{lem:bg}
Let $n\ge3$ and $L\ge1$, and let $f$ be a proper coloring of $K_n^{\times L}$ with $2n-2$ colors. Then there are a coordinate $i\in[L]$, a set $S$ of $n-2$ colors and a function $\gamma$ on $[n]$ such that $f(x)=\gamma(x_i)$ for every $x$ with $f(x)\notin S$.
\end{lemma}

Following the term ``dictator'' of~\cite{BG}, we call a coordinate $i$ as in Lemma~\ref{lem:bg} a \emph{dictator coordinate} of $f$. For $L=2$ we use the terms of~\cite[Section~2.4]{Z}. The vertices of $K_n^{\times2}$ are \emph{cells}, the cells with a fixed first coordinate form a \emph{row}, and those with a fixed second coordinate form a \emph{column}. Two cells of the same color agree in a coordinate. A color class that contains two cells of one row lies in that row, because a further cell of the class agrees with each of the two in a coordinate; it is called a \emph{row class}. \emph{Column classes} are defined in the same way. No class is of both kinds, and a class of neither kind is a single cell. A proper coloring of $K_n^{\times2}$ \emph{reads the first coordinate} if some column consists of cells of row classes, and \emph{reads the second coordinate} if some row consists of cells of column classes.

\begin{lemma}\label{lem:dictator}
Let $n\ge3$, and let $f,g$ be proper colorings of $K_n^{\times L}$ with colors from one set of $2n-2$ colors.
\begin{enumerate}
\item[(a)] The coloring $f$ has exactly one dictator coordinate, denoted $\delta(f)$. For $L=2$ it reads exactly one coordinate, namely the coordinate $\delta(f)$.
\item[(b)] Let $\sigma:[L]\to[L']$ be a map and put $f_\sigma(x)=f(x_{\sigma(1)},\ldots,x_{\sigma(L)})$ for $x\in[n]^{L'}$. Then $f_\sigma$ is a proper coloring of $K_n^{\times L'}$ and $\delta(f_\sigma)=\sigma(\delta(f))$.
\item[(c)] If $f(x)\ne g(y)$ whenever $x$ and $y$ differ in every coordinate, then $\delta(f)=\delta(g)$.
\end{enumerate}
\end{lemma}
\begin{proof}
\emph{Identifying coordinates.} In the situation of~(b), tuples $x,y\in[n]^{L'}$ that differ in every coordinate give tuples $(x_{\sigma(k)})_k$ and $(y_{\sigma(k)})_k$ that differ in every coordinate, so $f_\sigma$ is proper. If $i$ is a dictator coordinate of $f$ with $S$ and $\gamma$, then $f_\sigma(x)\notin S$ implies $f_\sigma(x)=\gamma(x_{\sigma(i)})$, so $\sigma(i)$ is a dictator coordinate of $f_\sigma$.

\emph{Two coordinates: exactly one coordinate is read.} Let $L=2$. If $f$ reads the first coordinate, the $n$ cells of some column lie in row classes, which are distinct because they lie in distinct rows; so $f$ has at least $n$ row classes. If it reads the second coordinate it has at least $n$ column classes. It cannot read both, since this would need $2n$ colors. Suppose that it reads neither. Then every column contains a cell outside the row classes, whose class is a column class or a single cell and lies in that column; likewise every row contains a row class or a class that is a single cell. Let $p$ be the number of rows that contain no row class and $q$ the number of columns that contain no column class. If $p=0$, there are at least $n$ row classes and at least one further class inside each column, hence at least $2n$ classes; the case $q=0$ is symmetric. If $p,q\ge1$, each of the $pq$ cells in these rows and columns is a class by itself, and there are at least $n-p$ row classes and $n-q$ column classes, hence at least
\[
 pq+(n-p)+(n-q)=2n-1+(p-1)(q-1)\ \ge\ 2n-1
\]
classes. Each case contradicts the number $2n-2$ of colors.

\emph{Two coordinates: the dictator coordinate.} Let $L=2$ and let the first coordinate be a dictator coordinate of $f$, with $S$ and $\gamma$. In the row with first coordinate $a$ every cell has its color in $S\cup\{\gamma(a)\}$, so a row has at most $n-1$ colors. A row that consists of cells of column classes has $n$ colors, because these classes lie in distinct columns. Hence $f$ does not read the second coordinate, and by the preceding paragraph it reads the first. In the same way, if the second coordinate is a dictator coordinate, $f$ reads the second coordinate. Since $f$ reads only one coordinate and has a dictator coordinate by Lemma~\ref{lem:bg}, it has exactly one, which is the coordinate that it reads.

\emph{Proof of~\textup{(a)} and~\textup{(b)}.} Suppose that $i\ne i'$ are dictator coordinates of a coloring $f$ of $K_n^{\times L}$. Let $\sigma:[L]\to[2]$ map $i$ to $1$ and every other coordinate to $2$. By the first paragraph both coordinates of $f_\sigma$ are dictator coordinates, which we have excluded. With Lemma~\ref{lem:bg} this proves~(a). Part~(b) follows from~(a) and the first paragraph.

\emph{Proof of~\textup{(c)}.} Suppose that $\delta(f)\ne\delta(g)$, and let $\sigma:[L]\to[2]$ map $\delta(f)$ to $1$ and every other coordinate to $2$. By~(b), $\delta(f_\sigma)=1$ and $\delta(g_\sigma)=2$, so $f_\sigma$ reads the first coordinate and $g_\sigma$ the second, and $f_\sigma(u)\ne g_\sigma(v)$ whenever the cells $u,v$ differ in both coordinates. Choose a column whose cells lie in row classes of $f_\sigma$, and let $t_a$ be the color of its cell in row $a$; these $n$ colors are distinct. If $g_\sigma(v)=t_a$ for a cell $v$ outside row $a$, then one of the at least two cells of the row class $f_\sigma^{-1}(t_a)$ differs from $v$ in both coordinates, which is excluded. So $g_\sigma$ uses $t_a$ only in row $a$, and no $t_a$ is the color of a column class of $g_\sigma$. As $g_\sigma$ has at least $n$ column classes, the two colorings together use at least $2n$ colors, a contradiction.
\end{proof}

\subsection{One substitution over a complete graph}\label{sec:kn-one}

Let $S'_\lambda(K_n)$ be the spanning subgraph of $S_\lambda(K_n)$ obtained by deleting every edge between two words of one basis chart that are not both constant on a common set $A_r$, $r\in\{1,2,3\}$. Edges at interior vertices of wires are kept, and so are, for $\lambda=0$, the edges between the two end grids of a wire. Two words of one end grid are constant on a common set $A_r$, so the end grids keep all their edges. Recall the physical graph $\mathcal P_n$ of~(Z1).

\begin{theorem}[One substitution over a complete graph]\label{thm:complete}
Let $n\ge3$ and $\lambda\ge0$. Then
\[
 \chi(S_\lambda(K_n))=2n-1\qquad\text{and}\qquad\cq(S_\lambda(K_n))=n .
\]
More generally, $\chi(G)=2n-1$ and $\cq(G)=n$ for every graph $G$ that admits graph homomorphisms $S'_\lambda(K_n)\to G\to\mathcal P_n$. This applies to $S'_\lambda(K_n)$, to $S_\lambda(K_n)$ and to $\mathcal P_n$.
\end{theorem}
\begin{proof}
Sending a vertex of $S_\lambda(K_n)$ to its measurement is a homomorphism $S_\lambda(K_n)\to\mathcal P_n$, because adjacent vertices carry measurements that are orthogonal at every label, hence distinct and adjacent. So the three named graphs are of the stated kind.

\emph{Upper bounds.} Let $S'_\lambda(K_n)\to G\to\mathcal P_n$. Then $\chi(G)\le\chi(\mathcal P_n)\le2n-1$ by~(Z1). The measurements that are the vertices of $\mathcal P_n$ form a quantum $n$-coloring of $\mathcal P_n$ in local dimension four, which pulls back to $G$; so $\cq(G)\le n$. The $n$ constant words of one basis chart are constant on every $A_r$ and form a clique of $S'_\lambda(K_n)$, whose image in $G$ is a clique on $n$ vertices; so $\cq(G)\ge n$ by~(Q4).

\emph{Lower bound.} It suffices to show that $S'_\lambda(K_n)$ has no proper coloring with $2n-2$ colors. Suppose that $f$ is one. For a basis $i$ and a set $D$ among $A_1,A_2,A_3,\{1\},\{2\},\{3\}$, let $f^i_D(x,y)$ be the color of the word of chart $i$ that takes $x$ on $D$ and $y$ on $D^c$. For fixed $i$ and $D$ these $n^2$ words are constant on a common set $A_r$: on $A_r$ if $D=A_r$, and on $A_r$ for either $r\ne s$ if $D=\{s\}$. Hence $f^i_D$ is a proper coloring of $K_n^{\times2}$ with $2n-2$ colors. Put $b^i_D=1$ if $\delta(f^i_D)=1$ and $b^i_D=0$ otherwise.

Fix a ternary template $\{A_r,\{s\},\{t\}\}$ at a basis $i$. The $n^3$ words of chart $i$ that are constant on $A_r$, indexed by their values on the three parts, induce a copy of $K_n^{\times3}$ in $S'_\lambda(K_n)$. Let $F$ be the restriction of $f$ to it. If $D$ is the $k$th part and $\sigma:[3]\to[2]$ maps $k$ to $1$ and the other two coordinates to $2$, then $f^i_D=F_\sigma$, so $\delta(f^i_D)=\sigma(\delta(F))$ by Lemma~\ref{lem:dictator}(b). Thus $b^i_D=1$ exactly for the part $D$ that is the coordinate $\delta(F)$, and
\[
 b^i_{A_r}+b^i_{\{s\}}+b^i_{\{t\}}=1 .
\]
Add the three equations of the three ternary templates at basis $i$ modulo two. Each singleton occurs in two of them, which leaves
\begin{equation}\label{eq:odd-bits}
 b^i_{A_1}+b^i_{A_2}+b^i_{A_3}\equiv1\pmod2\qquad(1\le i\le6).
\end{equation}

Now fix an edge $e=ik$ of $K_{3,3}$ and let $h_\ell(x,y)$ be the color of the vertex $(x,y)$ of layer $\ell$ of its wire. Each $h_\ell$ is a proper coloring of $K_n^{\times2}$, and $h_\ell(x,y)\ne h_{\ell+1}(x',y')$ whenever $x\ne x'$ and $y\ne y'$. By Lemma~\ref{lem:dictator}(c), $\delta(h_0)=\delta(h_{\lambda+1})$. By the identifications at the ends, $h_0=f^i_{A_{i,e}}$, and $h_{\lambda+1}=f^k_{A_{k,e}}$ if $e\ne36$, whereas $h_{\lambda+1}(x,y)=f^6_{A_3}(y,x)$ if $e=36$. Hence
\[
 b^i_{A_{i,e}}=b^k_{A_{k,e}}\quad(e\ne36),\qquad b^3_{A_3}=1-b^6_{A_3} .
\]
Each of the eighteen projections $P^i_r$ occurs once in~\eqref{eq:planes}, so these nine relations match the eighteen bits $b^i_{A_r}$ in pairs, eight pairs of equal bits and one pair of unequal bits. The sum of all eighteen bits is therefore odd. By~\eqref{eq:odd-bits} it is a sum of six odd numbers, hence even. This contradiction proves $\chi(S'_\lambda(K_n))\ge2n-1$, and $\chi(G)\ge2n-1$ follows for every $G$ as in the statement.
\end{proof}

\begin{remark}\label{rem:complete}
(i) For $n=3$ the equality $\chi(S_\lambda(K_3))=5$ is~\cite[Theorem~2.4]{Z}, whose proof uses~\cite[Lemma~2.3]{Z} in place of Lemma~\ref{lem:dictator} and no imported result. The proof above uses Lemma~\ref{lem:dictator} only for $L\le3$. For odd $n$ there is a second proof of the lower bound that does not use Lemma~\ref{lem:bg}: Theorem~\ref{thm:spheres} with $j=1$ gives $w^{2n-3}\ne0$ on $\B(S_\lambda(K_n))$, and also on $\B(S'_\lambda(K_n))$, because the map of~(T3) that its proof uses takes values there for every $n$, as the first step of the proof of Proposition~\ref{prop:even} shows.

(ii) Theorem~\ref{thm:complete} shows that the bound $\chi(S_\lambda(H))\le2\chi(H)-1$ of~(Z1) is attained by all complete graphs. It is also attained by every nonbipartite graph $H$ with $\chi(H)=3$: an odd cycle of $H$ gives a tower of length zero over $H$~\cite[proof of Theorem~7.1]{Z}, so $\chi(S_\lambda(H))\ge5$ by~(T4) and~(T1).

(iii) Ciardo proved that for every $n\ge3$ there are graphs with a quantum $n$-coloring in local dimension three and chromatic number at least $2n-1$~\cite[Corollary~35 and Example~36]{Ciardo24}. The graphs of Theorem~\ref{thm:complete} are explicit, have local dimension four, and their chromatic number is determined.
\end{remark}

\subsection{Substituting over a sphere}\label{sec:spheres}

The lower bound of~\cite{Z} starts from the circle $\B(K_3)$, which carries a tower of length zero. A sphere $S^{2h+1}$ with its standard circle action carries a tower of length $h$, so (T4) gives a tower of length $h+j$ after $j$ substitutions. We show that the first substitution does better: it produces a tower of length $2h+1$.

\paragraph{Presentations over a space.}
Let $Y$ be a compact polyhedron and put $X_0=Y$ and $X_g=S(X_{g-1})$ for $g\ge1$. The \emph{leaf points} of $X_g$ form the subspace $L_g$, where $L_0=Y$ and $L_g$ consists of the constant tuples $(z,z,z,z)$ in the block of basis $6$ with $z\in L_{g-1}$; thus $L_g$ is a copy of $Y$. A \emph{continuous presentation} of a point of $X_g$ is a tree as in Section~\ref{sec:substitution} with the following changes. A leaf carries a point of $Y$ and evaluates at level $g$ to the corresponding point of $L_g$. A basis node at level $g\ge1$ evaluates to the image of the values of its children under its template map. A wire node at level $g\ge1$ carries an edge $e$ of $K_{3,3}$ and a parameter $t\in[0,1]$ and evaluates to the point $(x,y,t)$ of the cylinder of $e$, where $x\in X_{g-1}$ is the value of its first child and $y\in L_{g-1}$ the value of its background, which is a leaf. Let $Z_{g,0}(Y)=L_g$, and for $m\ge1$ let $Z_{g,m}(Y)\subseteq X_g$ be the set of points that have a continuous presentation with at most $2m-1$ internal nodes and at most $m$ wire nodes. Each $Z_{g,m}(Y)$ is a finite union of images of products of copies of $Y$ and of intervals under piecewise linear maps, hence a compact polyhedron. A map $Y\to Y'$ induces maps $Z_{g,m}(Y)\to Z_{g,m}(Y')$. If $Y$ is a circle space, each $Z_{g,m}(Y)$ is invariant under the circle action and is a circle space. For $Y=S^1$ these are the spaces $Z_{g,m}$ of~\cite[Appendix~B.1]{Z}. Every point of $X_1=S(Y)$ has a presentation with one internal node, a basis node with three leaves or a wire node with two leaves, so
\begin{equation}\label{eq:first-level}
 Z_{1,m}(Y)=S(Y)\qquad(m\ge1).
\end{equation}

\begin{lemma}[Supports]\label{lem:supports}
Let $Y$ be a circle space and $g\ge1$. Let $x_0,\ldots,x_k$ be a tower on $X_g$ such that $x_0$ is a point of $L_g$ and $x_m\in C_{2m}(Z_{g,m}(Y);\Ftwo)$ for $1\le m\le k$. Then $X_{g+1}$ carries a tower $y_0,\ldots,y_{k+1}$ such that $y_0$ is a point of $L_{g+1}$ and $y_m\in C_{2m}(Z_{g+1,m}(Y);\Ftwo)$ for $1\le m\le k+1$.
\end{lemma}
\begin{proof}
Take the tower of~(T2). For $Y=S^1$ and $k=g$ this is the induction step in the proof of~\cite[Lemma~B.1(b)]{Z}, namely Step~4 of that proof. The argument uses neither assumption, and we recall it. Put $b=x_0$. In~\cite{Z} the chain $y_m$ is the coefficient of $u^m$ in a series $L+Jx$ in a formal variable $u$, where $x=\sum_mu^mx_m$. The series $L$ is the sum, over the nineteen templates, of the images under the template maps of \emph{corrected products} of copies of $x$, one for each part. A corrected product is a sum of terms, each obtained from a product of chains $x_{i_1},x_{i_2}$ and, for a ternary template, $x_{i_3}$ by letting chains of rotations act on the factors; a term has a \emph{correction order} $r\ge0$, where $r=0$ means that nothing is rotated, and contributes to the coefficient of $u^m$ with $m=r+\sum_\nu i_\nu$. Further, $Jx=Ex+u(BK+KB)x$, where $E$ inserts a chain on the positions $A_{b',e}$ at an end $b'$ of a wire with the point $b$ on the other positions, summed over the eighteen ends, and $K$ sends a chain $c$ to the chain swept by the points $(z,b,t)$, $z$ in $c$ and $t\in[0,1]$, in the nine cylinders. Rotations map $L_g$ and every $Z_{g,i}(Y)$ to themselves.

Fix $m$ with $1\le m\le k+1$. The terms of $L$ of correction order zero with exactly one factor of positive index sum to $Ex_m$ and cancel the term $Ex_m$ of $Jx$; for $m=k+1$ both are absent. Every other term of $y_m$ is of one of three kinds.
\begin{itemize}
\item A term of $L$ of correction order zero with at least two factors of positive indices $i_\nu$, which sum to $m$; its other factors are the point $b$. Each of its points has a presentation consisting of a basis node over presentations with at most $2i_\nu-1$ internal nodes and $i_\nu$ wire nodes and over leaves. As there are at least two indices $i_\nu$, it has at most $1+\sum_\nu(2i_\nu-1)\le2m-1$ internal nodes and at most $m$ wire nodes.
\item A term of $L$ of correction order $r\ge1$. The indices of its factors sum to $m-r$, a factor of index $i\ge1$ has its points in $Z_{g,i}(Y)$, and a factor of index zero is a point of $L_g$. Each of its points has a presentation with at most $1+2(m-r)\le2m-1$ internal nodes and at most $m-r$ wire nodes.
\item The terms $BKx_{m-1}$ and $KBx_{m-1}$. Their points are of the form $(z,y,t)$ with $z\in Z_{g,m-1}(Y)$ and $y$ in the orbit of $b$, which lies in $L_g$. Such a point is presented by a wire node over a presentation of $z$ and a leaf. For $m=1$ this is one internal node, which is a wire node; for $m\ge2$ there are at most $2(m-1)-1+1\le2m-1$ internal nodes and at most $m$ wire nodes.
\end{itemize}
Finally $y_0$ is the constant tuple on $b$ in the block of basis $6$, a point of $L_{g+1}$.
\end{proof}

\begin{theorem}[Maps from spheres]\label{thm:spheres}
Let $H$ be a graph, $h\ge0$ and $\varepsilon\in\{0,1\}$, and suppose that there is an equivariant map from the sphere $S^{2h+1+\varepsilon}$ with the antipodal involution to $\B(H)$. Let $j\ge1$ and $\lambda\ge0$, put $s=2h+j$, and let $G$ be $S^j_\lambda(H)$ or its induced subgraph $G^{(s)}_{j,\lambda}(H)$. Then
\[
 w^{2s+1+\varepsilon}\ne0\ \text{ on }\B(G)\qquad\text{and}\qquad\chi(G)\ \ge\ 2s+3+\varepsilon\;=\;4h+2j+3+\varepsilon .
\]
The graph $G^{(s)}_{j,\lambda}(H)$ contains a copy of $H$, has a $(\lambda+2)$-compatible quantum homomorphism to $H$ in local dimension $4^j$, and satisfies
\[
 |V|+|E|\ \le\ C^{\,s}\,(1+\lambda)^{\min\{s,\,3^{j-1}\}}
\]
with a constant $C$ that depends only on $|V(H)|$.
\end{theorem}

For comparison, (T4) alone gives $\chi(S^j_\lambda(H))\ge2(h+j)+3$ under the same hypothesis with $\varepsilon=0$.

\begin{proof}
Put $N=2h+1$ and let $Y=S^N$ be the unit sphere of $\C^{h+1}$ with the action of $S^1$ by scalar multiplication. It is a circle space whose half-turn is the antipodal map. Put $W=S(Y)$ and let $d:Y\to W$ send $x$ to the constant tuple $(x,x,x,x)$ in the block of basis $6$. Homology and cohomology have coefficients in $\Ftwo$.

\emph{Step 1: $d$ induces the zero map on $H_N$.} For a template with parts $D_1,\ldots,D_p$ and template map $f:Y^p\to W$, the composition of $f$ with the diagonal $Y\to Y^p$ is the map $d_i$ that sends $x$ to the constant tuple in the block of the basis $i$ of the template. By the K\"unneth formula, the diagonal $Y\to Y^p$ induces on $H_N$ the sum of the maps induced by the $p$ inclusions that put the argument into one factor and a fixed point $b$ into the others. Hence $(d_i)_*=\sum_D(a_D)_*$ on $H_N$, where $a_D:Y\to W$ inserts its argument on the part $D$ and $b$ on the other positions. All $d_i$ are homotopic to $d=d_6$: for an edge $e=ik$ of $K_{3,3}$ the points $(x,x,t)$ of the cylinder of $e$ give a homotopy from $d_i$ to $d_k$, and $K_{3,3}$ is connected. Summing over the nineteen templates therefore gives $19\,d_*=d_*$ on the left. On the right, every singleton part occurs in two ternary templates of its basis and cancels, and the parts $A_3$ of basis $6$ occur twice, in a ternary template and in the binary one. What remains are the insertions on $A_1,A_2,A_3$ at the bases $1,\ldots,5$ and on $A_1,A_2,A_3^c$ at basis $6$, that is, on $A_{b,e}$ for the eighteen ends $(b,e)$ of the nine wires. The two insertions at the ends of a wire $e$ are homotopic, through the points $(x,b,t)$ of its cylinder, so they cancel in pairs. Hence $d_*=0$ on $H_N(Y)$.

\emph{Step 2: if $N\ge3$, then $H_i(W)=0$ for $N+1<i<2N$.} A basis block is the space $\mathrm{Bl}(Y)$ of tuples in $Y^4$ with $x_0=x_r$ for some $r\ge1$. The projection $\mathrm{Bl}(Y)\to Y$ to $x_0$ is a fiber bundle: over an open set of $Y$ on which there is a continuous choice of rotations $\rho(x_0)\in SO(N+1)$ with $\rho(x_0)b=x_0$, the map $(x_0,x_1,x_2,x_3)\mapsto(x_0,\rho(x_0)^{-1}x_1,\rho(x_0)^{-1}x_2,\rho(x_0)^{-1}x_3)$ is a trivialization. The fiber is the set of triples in $Y^3$ with at least one entry equal to $b$. It is the union of the cells of dimensions $0$, $N$ and $2N$ in the product of the cell structures of $S^N$ with one vertex and one $N$-cell, so its homology is concentrated in the degrees $0,N,2N$. The base $S^N$ is simply connected, so the Serre spectral sequence~\cite{Serre} shows that $H_i(\mathrm{Bl}(Y))\ne0$ only for $i\in\{0,N,2N,3N\}$. Let $U\subset W$ be the union of the six blocks. By excision, $H_i(W,U)$ is the sum over the nine wires of $H_i(Y^2\times[0,1],Y^2\times\{0,1\})\cong H_{i-1}(Y^2)$, which is nonzero only for $i\in\{1,N+1,2N+1\}$. The exact sequence $H_i(U)\to H_i(W)\to H_i(W,U)$ gives the claim.

\emph{Step 3: a tower of length $2h+1$ on $W$.} Let $x_0$ be a point $b$ of $Y$. For $m=1,\ldots,h$ choose chains $x_m$ on $Y$ with $\partial x_m=Bx_{m-1}$. This is possible: $Bx_{m-1}$ is a cycle, since $\partial Bx_{m-1}=B\partial x_{m-1}$ vanishes for $m=1$ and equals $B^2x_{m-2}=0$ for $m\ge2$, and its dimension $2m-1$ satisfies $0<2m-1<N$. Put $x'_m=d_*x_m$ for $0\le m\le h$. As $d$ commutes with the circle actions, $d_*$ commutes with $B$, and $x'_0,\ldots,x'_h$ is a tower on $W$. The cycle $Bx'_h=d_*(Bx_h)$ has dimension $N$ and is a boundary by Step~1; choose $x'_{h+1}$ with $\partial x'_{h+1}=Bx'_h$. For $h+2\le m\le2h+1$, which occurs only if $N\ge3$, the cycle $Bx'_{m-1}$ has dimension $2m-1$ with $N+1<2m-1<2N$ and is a boundary by Step~2; choose $x'_m$ accordingly. This gives a tower $x'_0,\ldots,x'_{2h+1}$ on $W=X_1$ whose point $x'_0=d(b)$ lies in $L_1$.

\emph{Step 4: iteration.} By~\eqref{eq:first-level} the tower of Step~3 satisfies the hypothesis of Lemma~\ref{lem:supports} at level $g=1$ with $k=2h+1$. Applying that lemma $j-1$ times gives a tower $x_0,\ldots,x_s$ on $X_j=S^j(Y)$ with $x_0\in L_j$ and $x_m\in C_{2m}(Z_{j,m}(Y);\Ftwo)$. Since $L_j$ and $Z_{j,m}(Y)$, $m\le s$, lie in $Z_{j,s}(Y)$, it is a tower on the circle space $Z_{j,s}(Y)$ as well. In the rest of the proof $F$ stands for one of the two constructions $F(Y')=S^j(Y')$ and $F(Y')=Z_{j,s}(Y')$, applied to compact polyhedra $Y'$. By~(T1),
\begin{equation}\label{eq:odd-power}
 w_{F(Y)}^{\,2s+1}\ne0 .
\end{equation}

\emph{Step 5: one more power over the sphere of dimension $N+1$.} Let $\Sigma$ be the suspension of $Y$, the quotient of $Y\times[-1,1]$ in which $Y\times\{1\}$ and $Y\times\{-1\}$ are collapsed to two points. With the involution $[y,t]\mapsto[-y,-t]$ it is the sphere $S^{N+1}$ with the antipodal involution. For $t\in[-1,1]$ let $\varphi_t:Y\to\Sigma$ be the map $y\mapsto[y,t]$; the maps $\varphi_{\pm1}$ are constant. Put $M=F(Y)$ and $F_{\rm pt}=F(\mathrm{pt})$, let $p:M\to F_{\rm pt}$ be induced by $Y\to\mathrm{pt}$, and let $T$ be the half-turn of $M$; then $p\circ T=p$. Put $q=2s+1$. By~\eqref{eq:odd-power} there is a connected component $B_0$ of $M/T$ on which $w_M^q$ is nonzero. Let $M_0$ be its preimage in $M$. Let $P$ be the space obtained from $M_0\times[-1,1]$ and two copies $F_{\rm pt}^+,F_{\rm pt}^-$ of $F_{\rm pt}$ by identifying $(m,1)$ with $p(m)\in F_{\rm pt}^+$ and $(m,-1)$ with $p(m)\in F_{\rm pt}^-$. It is a compact polyhedron, and the involution $(m,t)\mapsto(Tm,-t)$, which exchanges $F_{\rm pt}^+$ and $F_{\rm pt}^-$, is free. Define
\[
 \Psi:P\to F(\Sigma),\qquad\Psi(m,t)=F(\varphi_t)(m)\ \text{ on }M_0\times[-1,1],
\]
and let $\Psi$ be induced on $F_{\rm pt}^+$ and $F_{\rm pt}^-$ by the inclusions of the two poles into $\Sigma$. This is well defined and continuous, because $\varphi_{\pm1}$ factors through a point, so that $F(\varphi_{\pm1})$ factors through $p$, and it is equivariant, because $\varphi_{-t}(-y)=-\varphi_t(y)$. It remains to show $w_P^{q+1}\ne0$.

The quotient $P/T$ contains a copy of $F_{\rm pt}$, the image of $F_{\rm pt}^+\cup F_{\rm pt}^-$, and $[m,t]\mapsto p(m)$ is a retraction of $P/T$ onto it. Hence $H^*(P/T)\to H^*(F_{\rm pt})$ is onto, and by the exact sequence of the pair the map $H^*(P/T,F_{\rm pt})\to H^*(P/T)$ is injective. Let $E_0$ be the quotient of $M_0\times[-1,1]$ by the relation $(m,t)\sim(Tm,-t)$. It is the unit disk bundle of the real line bundle over $B_0$ that belongs to the double cover $M_0\to B_0$, and its boundary $\partial E_0$, the image of $M_0\times\{\pm1\}$, is a copy of $M_0$. The space $P/T$ arises from $E_0$ and $F_{\rm pt}$ by attaching $\partial E_0$ to $F_{\rm pt}$ along the map induced by $p$. Hence collapsing $F_{\rm pt}$ in $P/T$ gives the same space as collapsing $\partial E_0$ in $E_0$. Both are pairs of a polyhedron and a subpolyhedron, so the map of pairs $(E_0,\partial E_0)\to(P/T,F_{\rm pt})$ induces an isomorphism $H^*(P/T,F_{\rm pt})\cong H^*(E_0,\partial E_0)$. Let $u\in H^1(E_0,\partial E_0)$ be the Thom class and $\pi:E_0\to B_0$ the projection. By the Thom isomorphism~\cite[\S\S8, 10]{MS}, $x\mapsto\pi^*x\cup u$ is an isomorphism $H^k(B_0)\to H^{k+1}(E_0,\partial E_0)$; in particular $H^1(E_0,\partial E_0)=\{0,u\}$, as $B_0$ is connected. The double cover $P\to P/T$ is trivial over $F_{\rm pt}$, so $w_P$ restricts to zero on $F_{\rm pt}$ and is the image of a class $\tilde u\in H^1(P/T,F_{\rm pt})$. Over the zero section $B_0\subset E_0$ the cover is $M_0\to B_0$, so $w_P$ restricts to $w_M|_{B_0}$ there, which is nonzero since its $q$th power is. Hence $\tilde u\ne0$, so $\tilde u$ corresponds to $u$, and the pullback of $w_P$ to $E_0$ is $\pi^*(w_M|_{B_0})$ because $\pi$ is a homotopy equivalence. Therefore $\tilde u\cup w_P^q\in H^{q+1}(P/T,F_{\rm pt})$ corresponds to $u\cup\pi^*(w_M^q|_{B_0})$, which is nonzero by the Thom isomorphism. Its image in $H^{q+1}(P/T)$ is $w_P^{q+1}$, which is nonzero by injectivity. As $\Psi$ pulls $w_{F(\Sigma)}$ back to $w_P$,
\begin{equation}\label{eq:even-power}
 w_{F(S^{N+1})}^{\,2s+2}\ne0 .
\end{equation}

\emph{Step 6: the graphs.} Let $\phi:S^{N+\varepsilon}\to\B(H)$ be the given equivariant map. Applying~(T3) $j$ times gives an equivariant map $S^j(S^{N+\varepsilon})\to\B(S^j_\lambda(H))$. We claim that it maps $Z_{j,s}(S^{N+\varepsilon})$ into $\B(G^{(s)}_{j,\lambda}(H))$. By the property stated in~(T3) and induction on the level, every vertex in a shore of the image of a point with a given continuous presentation has a presentation with the same tree: at a leaf carrying $y$ the vertex is a padded vertex of a shore of $\phi(y)$; at a basis node it is a word whose entries come from the shores of the children and are equal on each part; and at a wire node it is a pair in one of the at most two layers nearest to $(\lambda+1)t$ whose second entry comes from a leaf. So both shores consist of vertices of $G^{(s)}_{j,\lambda}(H)$. They are completely adjacent in $S^j_\lambda(H)$, hence in the induced subgraph. Thus $F(S^{N+\varepsilon})$ maps equivariantly to $\B(G)$ for both graphs $G$ of the statement, and~\eqref{eq:odd-power}, \eqref{eq:even-power} and~\eqref{eq:obstruction} give the two displayed assertions.

\emph{Step 7: the homomorphism and the order.} The padded vertices of $H$ are presented by single leaves, lie in $G^{(s)}_{j,\lambda}(H)$ and induce a copy of $H$. The $(\lambda+2)$-compatible quantum homomorphism $S^j_\lambda(H)\qto H$ of Section~\ref{sec:substitution} restricts to the induced subgraph, in which distances are not smaller. For the order we use the counting argument in the proof of~\cite[Theorem~3.6]{Z}, with the bounds $2s-1$ and $s$ for the numbers of internal nodes and of wire nodes and with the set $V(H)$ of labels at the leaves. A tree with at most $2s-1$ internal nodes with at most three children each has at most $6s-2$ nodes, so there are at most $C_0^{\,s}$ presentations in which the layers of the wire nodes are recorded only as $0$, $\lambda+1$ or interior, where $C_0$ depends only on $|V(H)|$. A presentation at level $j$ has at most $3^{j-1}$ wire nodes: at level one it has at most one, and if presentations at level $g-1\ge1$ have at most $c\ge1$ wire nodes, then a basis node at level $g$ has at most $3c$ below it and a wire node, whose background is a leaf, at most $1+c\le3c$ including itself. So the interior layers can be chosen in at most $(1+\lambda)^{\min\{s,3^{j-1}\}}$ ways, which bounds the number of vertices. For the edges, if $\lambda\le1$ there is at most one interior layer and nothing to choose. If $\lambda\ge2$, the argument of~\cite{Z} shows that once the layers of the presentation of one end of an edge are fixed, each interior layer in the presentation of the other end has at most three possible values. In both cases there are at most $C_0^{\,2s}3^s(1+\lambda)^{\min\{s,3^{j-1}\}}$ edges.
\end{proof}

\begin{corollary}[Iterated substitutions over complete graphs]\label{cor:kn}
Let $n\ge3$, $j\ge1$ and $\lambda\ge0$. Then
\[
 \chi(S^j_\lambda(K_n))\ \ge\
 \begin{cases}
 2n-3+2j,& n\text{ odd},\\
 2n-4+2j,& n\text{ even}.
 \end{cases}
\]
The same bound holds for the induced subgraph $G^{(s)}_{j,\lambda}(K_n)$ with $s=n-3+j$ for odd $n$ and $s=n-4+j$ for even $n$. This graph has quantum chromatic number $n$, a $(\lambda+2)$-compatible quantum $n$-coloring in local dimension $4^j$, and at most $C_n^{\,s}(1+\lambda)^s$ vertices and edges.
\end{corollary}
\begin{proof}
The box complex $\B(K_n)$ is the sphere $S^{n-2}$ with the antipodal involution. Apply Theorem~\ref{thm:spheres} to the identity map, with $2h+1=n-2$ and $\varepsilon=0$ for odd $n$, and with $2h+2=n-2$ and $\varepsilon=1$ for even $n$. Then $s=n-3+j$ and $2s+3=2n-3+2j$ in the first case, and $s=n-4+j$ and $2s+4=2n-4+2j$ in the second. The graph $G^{(s)}_{j,\lambda}(K_n)$ contains a copy of $K_n$ and has a quantum homomorphism to $K_n$, so its quantum chromatic number is $n$ by~(Q4).
\end{proof}

For $n=3$ the corollary is the bound $2j+3$ of~\cite{Z}. For odd $n$ and $j=1$ it gives $\chi(S_\lambda(K_n))\ge2n-1$, the exact value by Theorem~\ref{thm:complete}. For even $n$ and $j=1$ it gives $2n-2$, one less than the exact value. The next subsection shows that for $\lambda\ge2$ the method cannot give more.

\subsection{The height for even \texorpdfstring{$n$}{n}}\label{sec:even}

\begin{proposition}\label{prop:even}
Let $n\ge4$ be even and $\lambda\ge2$. On the box complexes of $S_\lambda(K_n)$ and of $S'_\lambda(K_n)$ the height of $w$ is $2n-4$. Thus~\eqref{eq:obstruction} certifies $2n-2$ colors for these graphs and no more, whereas their chromatic number is $2n-1$.
\end{proposition}
\begin{proof}
Put $N=n-2$, an even number, and identify $\B(K_n)$ with $S^N$ equivariantly. Here $N$ is the dimension of the sphere $\B(K_n)$, whereas in the proof of Theorem~\ref{thm:spheres} the letter stood for $2h+1$. Cohomology has coefficients in $\Ftwo$.

\emph{Lower bound.} Theorem~\ref{thm:spheres} with $H=K_n$, $2h+2=N$, $\varepsilon=1$ and $j=1$, so that $s=N-1$, gives $w^{2N}\ne0$ on $\B(S_\lambda(K_n))$. Its proof uses the map $S(S^N)\to\B(S_\lambda(K_n))$ of~(T3), and this map takes values in $\B(S'_\lambda(K_n))$, for even and for odd $n$: at a point $(x_0,x_1,x_2,x_3)$ of a basis block with $x_0=x_r$, all words in both shores are constant on $A_r$ by the property stated in~(T3), so opposite words are adjacent in $S'_\lambda(K_n)$, and at a point of a cylinder the adjacencies between the shores are among the edges that $S'_\lambda(K_n)$ keeps. Hence $w^{2N}\ne0$ on $\B(S'_\lambda(K_n))$ as well.

\emph{An auxiliary space.} Let $\mathcal W$ be the space formed by six \emph{full charts}, copies of $(S^N)^4$ indexed by the bases, and nine cylinders $(S^N)^2\times[0,1]$, attached as in $S(S^N)$: for $e=ik$ the point $(x,y,0)$ is the tuple of chart $i$ with $x$ on $A_{i,e}$ and $y$ on $A_{i,e}^c$, and $(x,y,1)$ is the tuple of chart $k$ with $x$ on $A_{k,e}$ and $y$ on $A_{k,e}^c$. The antipodal map in all coordinates is a free involution $T$ of $\mathcal W$. We construct an equivariant map $\Theta:\B(S_\lambda(K_n))\to\mathcal W$ and show $w^{2N+1}=0$ on $\mathcal W/T$. This proves $w^{2N+1}=0$ on $\B(S_\lambda(K_n))$, and on $\B(S'_\lambda(K_n))$ by restriction.

\emph{The map $\Theta$.} Assign to every vertex of $S_\lambda(K_n)$ a \emph{side}, which is a basis: a word of chart $b$ has side $b$, and an interior vertex of the wire of $e=ik$ has side $i$ if it lies in layer $1$ and side $k$ if it lies in a layer $\ell\ge2$. A vertex with side $b$ has four \emph{labels}: the entries of the word, or, for an interior vertex $(e,\ell,x,y)$, the label $x$ at the positions of $A_{b,e}$ and $y$ at the other two. Since $\lambda\ge2$, layers $1$ and $2$ are interior, and the definition of the edges shows two facts. If two adjacent vertices have the same side, their labels differ at all four positions. If they have different sides, they are interior vertices of one wire in layers $1$ and $2$.

Let $(\mu,\nu)\in\B(S_\lambda(K_n))$. If all vertices of both shores have the same side $b$, let $\mu_s$ and $\nu_s$ be the distributions of the label at position $s$ under $\mu$ and $\nu$. By the first fact $(\mu_s,\nu_s)\in\B(K_n)=S^N$, and we let $\Theta(\mu,\nu)$ be the point $((\mu_s,\nu_s))_{0\le s\le3}$ of the full chart of $b$. Otherwise some vertex $u$ of one shore and some vertex $v$ of the other have different sides, because two vertices of one shore with different sides cannot both have the side of a vertex of the other shore. By the second fact $u$ and $v$ lie in layers $1$ and $2$ of a wire $e$. Every vertex of the shore of $u$ is adjacent to the interior vertex $v$ and every vertex of the shore of $v$ is adjacent to $u$, so all vertices of both shores are pairs $(x,y)$ in layers of the wire of $e$. Let $(\mu_x,\nu_x)$ and $(\mu_y,\nu_y)$ be the distributions of the two entries, which are points of $\B(K_n)$, and let $t$ be the mean of the masses that $\mu$ and $\nu$ give to the layers $\ell\ge2$. We let $\Theta(\mu,\nu)$ be the point $((\mu_x,\nu_x),(\mu_y,\nu_y),t)$ of the cylinder of $e$; here $0<t<1$.

The map $\Theta$ commutes with the involutions. It is continuous on each cell of $\B(S_\lambda(K_n))$, hence continuous: on a cell all of whose vertices have one side this is clear, and on a cell with vertices of two sides, which lie on one wire $e$, the formula $((\mu_x,\nu_x),(\mu_y,\nu_y),t)$ is continuous on the whole cell and agrees with $\Theta$, because $t\in\{0,1\}$ exactly at the points whose shores have a single side, and there the point $(\cdot,\cdot,t)$ of the cylinder is the tuple of the chart at that end that repeats the two entries on $A_{b,e}$ and $A_{b,e}^c$.

\emph{The class $w^{2N+1}$ vanishes on $\mathcal W/T$.} Let $\mathcal A=(S^N)^4/T$ and $\mathcal E=(S^N)^2/T$, closed manifolds of dimensions $4N$ and $2N$; $\mathcal E$ is connected. The image $\mathcal R\subset\mathcal E$ of the antidiagonal $\{(x,-x)\}$ is a submanifold diffeomorphic to $\RP^N$ and disjoint from the image of the diagonal. The normal bundle of the antidiagonal in $(S^N)^2$ is the tangent bundle of $S^N$, by $v\mapsto(v,v)$, and $T$ acts on it as the differential of the antipodal map; so the normal bundle of $\mathcal R$ is isomorphic to the tangent bundle of $\RP^N$. Let $u\in H^N(\mathcal E)$ be the class dual to $\mathcal R$. Its restriction to $\mathcal R$ is the top Stiefel--Whitney class of the normal bundle~\cite[\S11]{MS}, so
\[
 \langle u^2,[\mathcal E]\rangle=\langle w_N(\RP^N),[\RP^N]\rangle\equiv\chi(\RP^N)=1\pmod2 ,
\]
as $N$ is even. Thus $u^2$ generates $H^{2N}(\mathcal E)=\Ftwo$. Moreover $u$ restricts to zero on the complement of $\mathcal R$, so $\beta^*u=0$ for the map $\beta:\mathcal E\to\mathcal E$, $[x,y]\mapsto[x,x]$, and $\tau^*u=u$ for the exchange $\tau[x,y]=[y,x]$, which maps $\mathcal R$ to itself.

Let $\mathcal A_b$ be the quotient of the full chart of $b$, let $p_{kl}:\mathcal A_b\to\mathcal E$ be the projection to the coordinates $k,l$, and for $\{r,a,c\}=\{1,2,3\}$ put
\[
 \kappa^b_r=p_{0a}^*u\cup p_{0c}^*u\in H^{2N}(\mathcal A_b).
\]
For an end $(b,e)$ of a wire let $\iota_{b,e}:\mathcal E\to\mathcal A_b$ be the attaching map, which puts $x$ on $A_{b,e}$ and $y$ on $A_{b,e}^c$. If $A_{b,e}=A_r$, then $p_{0a}\circ\iota_{b,e}$ and $p_{0c}\circ\iota_{b,e}$ are the identity and $\iota_{b,e}^*\kappa^b_r=u^2$. If $A_{b,e}=A_{r'}$ with $r'\ne r$, then one of the two compositions is $\beta$, because $r'\in\{a,c\}$, and $\iota_{b,e}^*\kappa^b_r=0$. At the end $(6,36)$, where $A_{6,36}=A_3^c$, both compositions are $\tau$ for $r=3$, so $\iota_{6,36}^*\kappa^6_3=u^2$, and for $r\in\{1,2\}$ the composition $p_{03}\circ\iota_{6,36}$ is $\beta\circ\tau$, so $\iota_{6,36}^*\kappa^6_r=0$. Consequently, for every end $(b,e)$ there is a class on $\mathcal A_b$ that restricts to the generator $u^2$ at this end and to zero at the other two ends at $b$.

Let $\mathcal U$ be the union of the six quotients $\mathcal A_b$ in $\mathcal W/T$. By excision, $H^{2N+1}(\mathcal W/T,\mathcal U)$ is the sum over the nine wires of $H^{2N+1}(\mathcal E\times[0,1],\mathcal E\times\{0,1\})\cong H^{2N}(\mathcal E)$, and the connecting map $H^{2N}(\mathcal U)\to H^{2N+1}(\mathcal W/T,\mathcal U)$ sends a class to the family of the sums of its restrictions to the two ends of each wire. By the preceding paragraph it is onto. The exact sequence of the pair then shows that the restriction $H^{2N+1}(\mathcal W/T)\to H^{2N+1}(\mathcal U)$ is injective. On $\mathcal A_b$ the class $w$ is pulled back from $\RP^N$ by the projection to the first coordinate, so $w^{N+1}$ vanishes there. Hence $w^{2N+1}$ restricts to zero on $\mathcal U$ and vanishes on $\mathcal W/T$.
\end{proof}

\begin{remark}\label{rem:short-wires}
The proof of the upper bound uses $\lambda\ge2$: the two layers on either side of the cut must be interior. Sending layer $\ell\ge1$ of every wire to layer $\ell-1$ and every word to itself is a homomorphism $S_{\lambda+1}(K_n)\to S_\lambda(K_n)$, so the height of $w$ does not increase with $\lambda$. For even $n$ and $\lambda\in\{0,1\}$ we only know that the height is $2n-4$ or $2n-3$: it is at least $2n-4$ by Theorem~\ref{thm:spheres} and at most $2n-3$ by~\eqref{eq:obstruction}, because the chromatic number is $2n-1$.
\end{remark}

\section{No equivariant map from the two-sphere}\label{sec:coindex}

The lower bounds for the chromatic number in~\cite{Z} and in Section~\ref{sec:complete} show that the box complex of a graph admits no equivariant map \emph{to} a sphere of low dimension. For graphs with a quantum $3$-coloring, maps in the other direction are restricted as well. The \emph{coindex} of a space $X$ with a free involution is the largest $k$ for which there is an equivariant map from the sphere $S^k$ with the antipodal involution to $X$~\cite{Mat}.

\begin{proposition}\label{prop:coindex}
Let $G$ be a graph with a quantum $3$-coloring. Then there is no equivariant map from the sphere $S^2$ with the antipodal involution to $\B(G)$. If $G$ is not bipartite, the coindex of $\B(G)$ is one.
\end{proposition}

The proof imports a theorem of Simons, Tardif and Wehlau, which we state first. Let $H$ be a graph. An \emph{arc} of $H$ is an ordered pair $(a,b)$ of adjacent vertices. Let $\Gamma(H)$ be the group generated by the arcs of $H$ with the relations
\begin{equation}\label{eq:stw-relation}
 (a,b)(c,b)^{-1}(c,d)(a,d)^{-1}=1\qquad\text{for all vertices with }a\sim b\sim c\sim d\sim a .
\end{equation}
If $a=c$ or $b=d$ the relation holds in the free group, so it does not matter whether such degenerate cycles are included. For a closed walk $v_0,v_1,\ldots,v_n=v_0$ of odd length $n$ in $H$, its \emph{signature} is the element
\[
 \prod_{i=0}^{n-1}(v_{2i},v_{2i+1})\,(v_{2i+2},v_{2i+1})^{-1}
\]
of $\Gamma(H)$, with the indices read modulo $n$ and the factors in the order of increasing $i$. In~\cite{STW} the box complex is realized as the order complex of the poset of pairs $(A,B)$ of nonempty sets of vertices with every vertex of $A$ adjacent to every vertex of $B$, with the involution $(A,B)\mapsto(B,A)$. The cells of $\B(H)$ in Section~\ref{sec:box} are indexed by this poset, so the order complex is the barycentric subdivision of $\B(H)$, and the two spaces are equivariantly homeomorphic.

\begin{theorem}[Simons, Tardif and Wehlau {\cite[Theorem~1 with Lemma~4, arXiv version~2]{STW}}]\label{thm:stw}
If there is an equivariant map from $S^2$ to $\B(H)$, then $H$ has a closed walk of odd length whose signature is the identity of $\Gamma(H)$.
\end{theorem}

\begin{proof}[Proof of Proposition~\ref{prop:coindex}]
Index the colors by $\mathbb Z/3$ and let $(P_{u,a})$ be a quantum $3$-coloring of $G$ on $\C^d$. By~(Q1) the projectors at adjacent vertices commute. For an arc $(u,v)$ and $\theta\in\R$ put
\[
 U_\theta(u,v)=\sum_{a\in\mathbb Z/3}\bigl(e^{i\theta}P_{u,a}P_{v,a+1}+e^{-i\theta}P_{u,a}P_{v,a-1}\bigr).
\]
The nine products $P_{u,a}P_{v,b}$ are pairwise orthogonal projections with sum $I$, and the three with $a=b$ vanish by the edge rule. So $U_\theta(u,v)$ is unitary, and $U_\theta(v,u)=U_\theta(u,v)^{-1}$, because $P_{u,a}P_{v,a+1}$ has the phase $e^{i\theta}$ in $U_\theta(u,v)$ and $e^{-i\theta}$ in $U_\theta(v,u)$.

\emph{The product along two edges depends only on the ends.} Let $u\sim v\sim w$, where $u=w$ is allowed. We claim that
\begin{equation}\label{eq:two-edges}
 U_\theta(u,v)\,U_\theta(v,w)=\sum_{a\in\mathbb Z/3}P_{u,a}\bigl(P_{w,a}+e^{-2i\theta}P_{w,a+1}+e^{2i\theta}P_{w,a-1}\bigr).
\end{equation}
Indeed, the left side is the sum of the terms $\varphi(b-a)\varphi(c-b)\,P_{u,a}P_{v,b}P_{w,c}$ over all $a,b,c$ with $b\ne a$ and $c\ne b$, where $\varphi(1)=e^{i\theta}$ and $\varphi(-1)=e^{-i\theta}$. If $c=a$, the two values $b=a\pm1$ occur, both with phase $1$, and $P_{u,a}(P_{v,a+1}+P_{v,a-1})P_{w,a}=P_{u,a}(I-P_{v,a})P_{w,a}=P_{u,a}P_{w,a}$, because $P_{u,a}P_{v,a}=0$. If $c=a+1$, only $b=a-1$ occurs, with phase $e^{-2i\theta}$, and $P_{u,a}P_{v,a-1}P_{w,a+1}=P_{u,a}(I-P_{v,a}-P_{v,a+1})P_{w,a+1}=P_{u,a}P_{w,a+1}$, because $P_{u,a}P_{v,a}=0$ and $P_{v,a+1}P_{w,a+1}=0$. The case $c=a-1$ is symmetric, with phase $e^{2i\theta}$.

\emph{A representation of $\Gamma(G)$.} By~\eqref{eq:two-edges}, for $a\sim b\sim c\sim d\sim a$ the products $U_\theta(a,b)U_\theta(b,c)$ and $U_\theta(a,d)U_\theta(d,c)$ are equal. As $U_\theta(c,b)^{-1}=U_\theta(b,c)$ and $U_\theta(a,d)^{-1}=U_\theta(d,a)$, the assignment $(u,v)\mapsto U_\theta(u,v)$ respects the relations~\eqref{eq:stw-relation} and defines a homomorphism from $\Gamma(G)$ to the unitary group of $\C^d$.

\emph{Conclusion.} Suppose that there is an equivariant map $S^2\to\B(G)$. By Theorem~\ref{thm:stw} there is a closed walk $v_0,\ldots,v_n=v_0$ of odd length with trivial signature. The homomorphism sends the factor $(v_{2i},v_{2i+1})(v_{2i+2},v_{2i+1})^{-1}$ of the signature to $U_\theta(v_{2i},v_{2i+1})U_\theta(v_{2i+1},v_{2i+2})$. As $i$ runs from $0$ to $n-1$ and the indices are read modulo the odd number $n$, the walk is traversed twice. So the signature is sent to $H(\theta)^2$, where
\[
 H(\theta)=U_\theta(v_0,v_1)\,U_\theta(v_1,v_2)\cdots U_\theta(v_{n-1},v_n),
\]
and $H(\theta)^2=I$ for every $\theta$. Thus $H(\theta)$ is a unitary involution, $(I-H(\theta))/2$ is a projection that depends continuously on $\theta$, and its rank is constant. For $\theta=2\pi/3$ put $Z_u=\sum_ae^{2\pi ia/3}P_{u,a}$. Then $U_\theta(u,v)=Z_u^{-1}Z_v$ on every arc, so the product telescopes and $H(2\pi/3)=Z_{v_0}^{-1}Z_{v_n}=I$. For $\theta=\pi$ every $U_\theta(u,v)$ equals $-I$, so $H(\pi)=(-1)^nI=-I$. The rank of $(I-H(\theta))/2$ is $0$ in the first case and $d$ in the second, a contradiction.

Finally, if $G$ contains an odd cycle $C$, then $\B(C)$ is a circle on which the exchange of the shores is the half-turn~\cite[proof of Theorem~7.1]{Z}, and the inclusion $\B(C)\subseteq\B(G)$ is an equivariant map from $S^1$. So the coindex is at least one, and it is at most one by the first part.
\end{proof}

The representation $U_\theta$ interpolates between two homomorphisms that every quantum $3$-coloring provides, the telescoping one at $\theta=2\pi/3$ and the sign of the length at $\theta=\pi$. The proposition restricts the graphs with a quantum $3$-coloring from the side opposite to the chromatic number: their box complexes can have large height, as the graphs of~\cite{Z} show, but their coindex is at most one.

\section{Further small graphs}\label{sec:small}

\subsection{Chromatic number four at a prescribed compatibility parameter}\label{sec:stretch}

We recall the graph $F_{22}$ and its coloring from~\cite[Proposition~6.1]{Z}. It has a \emph{center} $o$ adjacent to nine \emph{terminals} $A,\bar A,B,\bar B,C,\bar C,X,Y,Z$, the three edges $A\bar A$, $B\bar B$, $C\bar C$, and six \emph{private edges} $u_rv_r$, where $u_r$ is adjacent to the two terminals of the left pair and $v_r$ to the two terminals of the right pair in row $r$ of the table
\[
\begin{array}{c|cc}
 r&\text{left}&\text{right}\\ \hline
 1&A,X&Z,\bar C\\
 2&\bar A,X&C,Z\\
 3&X,B&Y,\bar C
\end{array}
\qquad\qquad
\begin{array}{c|cc}
 r&\text{left}&\text{right}\\ \hline
 4&X,\bar B&Y,C\\
 5&A,Y&B,Z\\
 6&\bar A,Y&\bar B,Z
\end{array}
\]
With the Pauli matrices $x,y,z$, the unbarred terminals carry the two-qubit observables
\[
 A=I\otimes z,\quad B=x\otimes x,\quad C=z\otimes x,\quad X=x\otimes I,\quad Y=z\otimes I,\quad Z=y\otimes y,
\]
and a barred terminal $\bar S$ carries $-S$. The two observables of every pair in the table commute, and $AX=CZ$, $XB=YC$, $AY=-BZ$. The quantum $3$-coloring with colors $0,1,2$ is as follows. The center returns $0$. A terminal with observable $S$ returns $1$ on the eigenspace of $S$ for $+1$ and $2$ on the eigenspace for $-1$. An end of a private edge whose pair carries the observables $U,V$ returns $0,1,2$ on the ranges of
\begin{equation}\label{eq:private}
 \tfrac12(I-UV),\qquad\tfrac14(I-U)(I-V),\qquad\tfrac14(I+U)(I+V).
\end{equation}
By the three relations, the products $UV$ at the two ends of a private edge are negatives of each other.

\begin{proposition}[Stretching $F_{22}$]\label{prop:stretch}
Let $R\ge2$ and $l=\lceil(R-1)/2\rceil$. There is a graph $F_{22}^{(l)}$ with $22+69\,l$ vertices and $42+115\,l$ edges, chromatic number four and quantum chromatic number three, that has an $R$-compatible quantum $3$-coloring in local dimension four. Every quantum $3$-coloring of it has local dimension divisible by four. In particular $F_{22}^{(1)}$, with $91$ vertices and $157$ edges, serves $R=2$ and $R=3$.
\end{proposition}
\begin{proof}
\emph{Contexts.} We call a set of vertices of $F_{22}$ whose measurements commute pairwise a \emph{context}, the usual word for a set of jointly measurable observables. Up to bars, the pairs of the table are $AX$, $XB$, $BZ$, $ZC$, $CY$, $YA$, each occurring in two rows. We use twelve contexts. For each of the six pairs take its two terminals, the barred partner of its terminal in $\{A,B,C\}$, and the two ends of private edges whose pairs consist of these terminals; to the contexts of $AX$, $BZ$ and $CY$ add the center. The other six contexts are the pairs $\{u_r,v_r\}$. The measurements of one of the first six contexts are functions of two commuting observables, and the measurement of $o$ is trivial, so they commute. For a private edge with product $M=UV$ at $u_r$, the projections of $u_r$ are $\frac12(I-M)$ and two projections below $\frac12(I+M)$ by~\eqref{eq:private}, and those of $v_r$ are $\frac12(I+M)$ and two projections below $\frac12(I-M)$; any projection of the first kind and any of the second are orthogonal or comparable, hence commute. Every edge of $F_{22}$ lies in a context: the contexts of $AX$, $BZ$, $CY$ contain $o$ and together all nine terminals, the edge $A\bar A$ lies in the context of $AX$, and likewise $B\bar B$ and $C\bar C$; an edge from an end of a private edge to a terminal lies in the context of its pair; and a private edge is a context. The center lies in three contexts and every other vertex in two.

\emph{The graph.} For every context $\mathfrak c$ and every vertex $x\in\mathfrak c$ take a vertex $(x,\mathfrak c)$, a \emph{copy} of $x$; there are $3+2\cdot21=45$ copies. For every edge $xy$ of $F_{22}$ choose one context that contains it and join the copies of $x$ and $y$ in that context. A \emph{diamond} is the graph $K_4$ with one edge removed, and the two ends of the removed edge are its \emph{tips}. A \emph{chain} of $l$ diamonds consists of $l$ diamonds in a row, consecutive ones sharing a tip; its two \emph{ends} are the tips at its two extremities. It has $3l+1$ vertices and $5l$ edges, and its ends are at distance $2l$. For every vertex $x$ of $F_{22}$ join the copies of $x$ by chains of $l$ diamonds, two chains in a row for the three copies of $o$ and one chain for the two copies of any other vertex, the copies being the ends of the chains. This uses $23$ chains and gives a graph $F_{22}^{(l)}$ with $45+23(3l-1)=22+69\,l$ vertices and $42+23\cdot5l$ edges.

\emph{The coloring.} Every copy of $x$ and every tip on a chain between copies of $x$ carries the measurement $(P_0,P_1,P_2)$ of $x$, and the two other vertices of each diamond carry $(P_1,P_2,P_0)$ and $(P_2,P_0,P_1)$. The edge rule holds between copies as in $F_{22}$, and on the edges of a diamond because projectors for different outcomes of one measurement are orthogonal. The local dimension is four. A diamond contains a triangle, so the quantum chromatic number is three by~(Q3).

\emph{Chromatic number.} In a proper $3$-coloring of a diamond the two tips have the same color. A proper $3$-coloring of $F_{22}^{(l)}$ would thus give all copies of each vertex $x$ of $F_{22}$ the same color, and these colors would form a proper $3$-coloring of $F_{22}$, which does not exist. A proper $4$-coloring of $F_{22}$ extends: copies and tips receive the color of their vertex, and the two other vertices of a diamond two of the remaining colors.

\emph{Compatibility.} Since $2l+1\ge R$, it suffices to show that the measurements at the two ends of a path $\gamma$ with at most $2l+1$ edges commute. All vertices of a chain between copies of $x$ carry the measurement of $x$ up to a permutation of the outcomes, and every edge of the graph lies in a chain or joins two copies in one context. If the copies on $\gamma$ lie in at most one context $\mathfrak c$, then every vertex of $\gamma$ is a copy in $\mathfrak c$ or lies on a chain with an end in $\mathfrak c$, or $\gamma$ lies in a single chain; in each case the measurements along $\gamma$ are, up to permutations of outcomes, measurements of one context, or of one vertex, and commute. Otherwise $\gamma$ contains a subpath that leads from a copy in one context to a copy in another one with no copy in between. This subpath runs through a whole chain, between copies of some vertex $x$, and has at least $2l$ edges. So $\gamma$ has at most one further edge. One end of $\gamma$ is then an end of this chain and carries the measurement of $x$. The other is the other end of the chain, or a neighbor of it, which is a copy in the same context or a vertex of a chain between copies of $x$; its measurement commutes with that of $x$.

\emph{Divisibility.} Consider any quantum $3$-coloring of $F_{22}^{(l)}$. A diamond is the union of two triangles with a common edge, so by~(Q2) its two tips carry the same projectors. Hence all copies of a vertex $x$ of $F_{22}$ carry the same measurement, and assigning it to $x$ gives a quantum $3$-coloring of $F_{22}$ in the same local dimension, which is divisible by four by~(Z6).
\end{proof}

For $l=1,2$ the program \texttt{check\_note.py} finds that the measurements of this coloring commute up to distance $2l+1$ and not beyond.

\subsection{The 165 measurements with fewer edges}\label{sec:165}

We recall the graph $F_{\rm phys}$ and the tower over it from~\cite[Section~6.2 and Appendix~E]{Z}. Let $C=\mathbb Z/3$ be the set of colors, and index the positions $1,2,3$ of a basis by $i\in C$ in this order. Thirteen \emph{difference triples} in $C^3$ are used:
\[
 O=(0,0,0),\qquad\alpha_i=\mathbf1-e_i,\qquad\beta_i=-\alpha_i,\qquad\delta_i=-e_i,\qquad\upsilon_i=(s-i)_{s\in C}\qquad(i\in C).
\]
For $k\in C$ and a difference triple $p$ the word $\omega(k,p)=(k,k+p_0,k+p_1,k+p_2)$ is a word of the basis chart of $S_0(K_3)$ at each basis. Its measurement in basis $b$ is written $p/k$ when the basis is clear. The graph $F_{\rm phys}$ is the physical graph on the measurements of these $39$ words in the six bases. It has $165$ vertices: the three \emph{scalar} measurements, from $p=O$, which occur in all six bases; six measurements on each of the nine planes, from $p=\alpha_i,\beta_i$, each occurring in the two bases of its plane; and $108$ measurements with an outcome of rank one, from $p=\delta_i,\upsilon_i$, each occurring in one basis. In basis $b$ and for each $i\in C$, the nine words $\omega(k,p)$ with $p\in\{O,\alpha_i,\beta_i\}$ take one color on the positions $\{0,i+1\}$ and one color on the other two; for the three values of $i$ these are the end grids of the three wires at $b$.

For $p\ne q$ in $C^3$ we say that $p$ \emph{precedes} $q$ if all coordinates of $q-p$ are $0$ or $1$. Let $\Delta$ be the simplicial complex on the thirteen difference triples whose simplices are the sets $\{p_0,\ldots,p_m\}$ in which $p_a$ precedes $p_{a'}$ whenever $a<a'$; it is a subcomplex of the triangulation of the torus $(\R/3\mathbb Z)^3$ by the ordered simplices of its unit cubes. Any two vertices of a simplex are comparable. It contains the $21$ triangles
\begin{equation}\label{eq:triangles}
\begin{gathered}
 (O,\delta_i,\alpha_i),\quad(\delta_i,\alpha_i,\upsilon_{i+1}),\quad(\alpha_i,\upsilon_{i+1},\upsilon_i),\quad(\alpha_i,\upsilon_i,\beta_i),\\
 (\beta_i,\upsilon_i,\upsilon_{i+2}),\quad(\beta_i,\upsilon_{i+2},\delta_{i+1}),\quad(\beta_i,\delta_{i+1},O)\qquad(i\in C).
\end{gathered}
\end{equation}
Let $\Delta_{21}$ be their union. It has $36$ edges: the nine edges $O\alpha_i$, $\alpha_i\beta_i$, $\beta_iO$, each in one triangle, and $27$ edges that lie in two triangles each. The sum of the $21$ triangles therefore has the three loops $O,\alpha_i,\beta_i,O$ as its boundary modulo two.

The circle space $Z$ of~\cite[Appendix~E]{Z} consists of six copies of $(\R/3\mathbb Z)\times\Delta$, on which the circle $\R/3\mathbb Z$ acts by translation of the first coordinate $t$, and of nine cylinders that join, for each edge $e$ of $K_{3,3}$, the tori over the loops of the two end grids of $e$, with a twist at the complemented edge. The map $Z\to\B(F_{\rm phys})$ sends the point $(t,p)$ of the copy of basis $b$ to the pair of the distributions of the measurement of the word
\[
 \bigl(\lfloor\sigma\rfloor,\lfloor\sigma+p_0\rfloor,\lfloor\sigma+p_1\rfloor,\lfloor\sigma+p_2\rfloor\bigr)\bmod3
\]
in basis $b$ for $\sigma$ uniform in $[t,t+\frac12]$ and for $\sigma$ uniform in $[t+\frac32,t+2]$, and it is constant along the cylinders. The words that occur at $(t,p)$ are words $\omega(k,p')$ with $p'$ a vertex of the smallest simplex of $\Delta$ that contains $p$, and two words in opposite shores differ at all four positions. The tower consists of a point $x_0$ with $p=O$ and a $2$-chain $x_1$ with $\partial x_1=Bx_0$, which exists because the sum of the three loops bounds in $\Delta$ and the cylinder of the complemented edge is attached with a twist.

Let $Z_{21}\subseteq Z$ be the subspace built in the same way from $\Delta_{21}$. It is invariant under the circle action and contains the loops, the cylinders and the $21$ triangles, and $\Delta_{21}$ is connected; the argument of~\cite[Appendix~E]{Z} for the existence of $x_1$ uses only this. So $Z_{21}$ is a circle space with a tower $x_0,x_1$ of length one.

\paragraph{Two subgraphs.}
Let $F_\triangle$ be the spanning subgraph of $F_{\rm phys}$ with the following edges, for each basis: for each difference triple $p$, the three edges among the measurements $p/0,p/1,p/2$; and for each edge $\{p,q\}$ of $\Delta_{21}$, with $p$ preceding $q$, the three edges from $p/k$ to $q/(k+1)$, $k\in C$. These pairs are edges of $F_{\rm phys}$, because the words $\omega(k,p)$ and $\omega(k+1,q)$ differ at all four positions. Let $F_*$ be obtained from $F_\triangle$ by deleting, for each basis and each $i\in C$, the three edges among $\delta_i/0,\delta_i/1,\delta_i/2$.

The following lemma allows the deletion of edges without loss of the tower.

\begin{lemma}[Deleting an edge]\label{lem:link}
Let $uv$ be an edge of a graph $G$, and let $\Lambda$ be the simplicial complex whose vertex set is the disjoint union of $N(v)\setminus\{u\}$, the \emph{left} vertices, and $N(u)\setminus\{v\}$, the \emph{right} vertices, and whose simplices are the sets in which every left vertex is adjacent in $G$ to every right vertex. If $\Lambda$ is contractible, there is an equivariant map $\B(G)\to\B(G-uv)$ that is the identity on $\B(G-uv)$. In particular, every tower over $G$ gives a tower over $G-uv$ on the same circle space.
\end{lemma}
\begin{proof}
A point $(\mu,\nu)$ of $\B(G)$ lies outside $\B(G-uv)$ exactly if $u$ and $v$ lie in opposite shores. Let $\Sigma_{uv}$ be the closure of the set of points with $\mu(u)>0$ and $\nu(v)>0$. Its points are described by the weights $\xi$ of $\mu$ off $u$ and $\eta$ of $\nu$ off $v$: these are nonnegative functions on the left and on the right vertices, respectively, with sums at most one, such that every left vertex in the support of $\xi$ is adjacent to every right vertex in the support of $\eta$. The point lies in $\B(G-uv)$ exactly if one of the two sums equals one. Hence $\Sigma_{uv}$ is the cone, with apex $\xi=\eta=0$, over $\Sigma_{uv}\cap\B(G-uv)$, the set where the larger of the two sums is one; and radial projection maps this set homeomorphically onto the set where the two sums add up to one, which is $\Lambda$. The corresponding set $\Sigma_{vu}$ is the image of $\Sigma_{uv}$ under the involution, and the two meet only in $\B(G-uv)$, because $u$ is not adjacent to itself. So $\B(G)$ arises from $\B(G-uv)$ by attaching two cones over copies of $\Lambda$, which the involution exchanges. If $\Lambda$ is contractible, the base of a cone is a deformation retract of the cone~\cite[Chapter~0]{Hatcher}. Choosing a retraction on $\Sigma_{uv}$ and its conjugate by the involution on $\Sigma_{vu}$ gives the required map. Composing with it turns an equivariant map from a circle space to $\B(G)$ into one to $\B(G-uv)$.
\end{proof}

\begin{theorem}\label{thm:star}
The graphs $F_\triangle$ and $F_*$ are spanning subgraphs of $F_{\rm phys}$ with $732$ and $678$ edges. There is a tower of length one over each of them. Consequently $w^3\ne0$ on their box complexes, $(\cq,\chi)=(3,5)$ for both, and each has a quantum $3$-coloring in local dimension four and none in a smaller local dimension.
\end{theorem}
\begin{proof}
\emph{Numbers of edges.} Fix a basis $b$. The definition of $F_\triangle$ lists $13\cdot3+36\cdot3=147$ pairs of words of $b$. They are distinct pairs, and distinct words of one basis have distinct measurements, so they give a set $E_b$ of $147$ edges. No edge of $\Delta_{21}$ joins two of the triples $\alpha_i,\beta_i$ with different indices, as the list~\eqref{eq:triangles} shows. Hence the edges of $E_b$ both of whose ends occur in further bases are the edges inside the three end grids. In the end grid with index $i$ the definition gives the three edges among the scalar measurements, the six edges among the $\alpha_i/k$ and among the $\beta_i/k$, and nine edges from the three edges $O\alpha_i$, $\alpha_i\beta_i$, $\beta_iO$ of $\Delta_{21}$. These are $18$ edges, and the grid $K_3^{\times2}$ has $18$ edges, so every end grid is complete in $E_b$. There are $3+3\cdot15=48$ such edges in $E_b$, and the other $99$ edges of $E_b$ have an end with an outcome of rank one, which belongs to basis $b$ only. Two bases that share a plane give the same $18$ edges on its grid. Altogether $F_\triangle$ has $3+9\cdot15+6\cdot99=732$ edges. The deleted edges join measurements that belong to one basis, so they are $6\cdot3\cdot3=54$ distinct edges, and $F_*$ has $678$ edges. The same count shows that the subgraph of $F_\triangle$ induced on the $39$ measurements of basis $b$ has edge set $E_b$.

\emph{A tower over $F_\triangle$.} We show that the map $Z\to\B(F_{\rm phys})$ sends $Z_{21}$ into $\B(F_\triangle)$. Let $(t,p)$ be a point of a copy of $(\R/3\mathbb Z)\times\Delta_{21}$, and let $\omega(k,p')$ and $\omega(k',q')$ be words in opposite shores of its image. Then $p'$ and $q'$ are vertices of a triangle of $\Delta_{21}$ and the two words differ at all four positions. If $p'=q'$, then $k\ne k'$ and the edge is one of the first kind. Otherwise we may assume that $p'$ precedes $q'$. Position $0$ gives $k'\ne k$. Some coordinate of $q'-p'$ is $1$, and the corresponding position gives $k'\ne k-1$. So $k'=k+1$ and the edge is of the second kind. The points of the cylinders have the same images as points of the copies. With the tower $x_0,x_1$ on $Z_{21}$ this is a tower of length one over $F_\triangle$.

\emph{A tower over $F_*$.} We delete the $54$ edges one by one and apply Lemma~\ref{lem:link} each time. Consider first the edge between $u=\delta_0/0$ and $v=\delta_0/1$ of a basis $b$ in $F_\triangle$. The edges of $\Delta_{21}$ at $\delta_0$ are $\delta_0O$, $\delta_0\alpha_0$, $\delta_0\upsilon_1$ and $\beta_2\delta_0$, each written with the preceding triple first. So the neighbors of $\delta_0/k$ are $\delta_0/(k\pm1)$, $O/(k+1)$, $\alpha_0/(k+1)$, $\upsilon_1/(k+1)$ and $\beta_2/(k-1)$, all in basis $b$. The left and right vertices of $\Lambda$ and their adjacencies in $F_\triangle$ are given by the following table, with the left vertices as rows and an entry $1$ for adjacent pairs.
\[
\begin{array}{c|ccccc}
  &\alpha_0/1&\beta_2/2&\delta_0/2&O/1&\upsilon_1/1\\ \hline
\alpha_0/2&1&0&0&1&1\\
\beta_2/0&0&1&0&1&1\\
\delta_0/2&0&0&0&0&0\\
O/2&0&0&0&1&0\\
\upsilon_1/2&0&0&0&0&1
\end{array}
\]
The entries follow from the definition of $E_b$: among the triples $O,\alpha_0,\beta_2,\delta_0,\upsilon_1$ the edges of $\Delta_{21}$ are, besides those at $\delta_0$, the edges $O\alpha_0$, $\beta_2O$, $\upsilon_1\alpha_0$ and $\beta_2\upsilon_1$, again with the preceding triple first.

If the right neighbors of a left vertex $a$ of $\Lambda$ are among those of another left vertex $a'$, then every simplex that contains $a$ remains a simplex when $a'$ is added, so the link of $a$ is a cone and removing $a$ does not change the homotopy type; the same holds with left and right exchanged. In the table, the last three rows are dominated in this sense by the first row. After their removal the first three columns are dominated by the fourth. The remaining two rows and two columns have all entries equal to $1$ and span a simplex. So $\Lambda$ is contractible.

The set $E_b$ is invariant under adding a constant to all colors and under cyclic permutations of the positions $1,2,3$, which permute the difference triples and the triangles~\eqref{eq:triangles} and preserve the relation of preceding. Hence for each of the $54$ edges the complex $\Lambda$ formed in $F_\triangle$ is given by the same table up to these symmetries and up to an exchange of left and right. When some of the $54$ edges have already been deleted, the only possible change of the complex $\Lambda$ of a further one is that the vertex in the third row, in the third column, or both, are absent, because the other vertices are not ends of deleted edges and all entries of that row and column are zero. The reduction above still applies. By Lemma~\ref{lem:link}, applied $54$ times, the tower over $F_\triangle$ gives a tower over $F_*$.

\emph{Chromatic numbers and local dimension.} By~(T1), $w^3\ne0$ on both box complexes and both chromatic numbers are at least five. They are at most five because $\chi(F_{\rm phys})=5$ by~(Z7). The measurements themselves form a quantum $3$-coloring in local dimension four, and the three scalar measurements form a triangle in $F_*$, so $\cq=3$ by~(Q3). Since $w^2\ne0$, every quantum $3$-coloring has even local dimension by~(Z5), and local dimension two is excluded by~(Z4).
\end{proof}

\subsection{Chromatic number five at a prescribed compatibility radius}\label{sec:FR}

Let $R\ge2$ and $\lambda=R-2$. The graph $F_R$ consists of six \emph{charts} and nine wires. The chart of basis $b$ has the $39$ words $\omega(k,p)$ of basis $b$ as vertices and the $147$ pairs of words that define $E_b$ as edges; the six charts are disjoint. For an edge $e=ik$ of $K_{3,3}$, layer $0$ of its wire is the end grid of $e$ in chart $i$ and layer $\lambda+1$ is the end grid of $e$ in chart $k$, where the cell $(x,y)$ is the word that takes $x$ on $A_{b,e}$ and $y$ on $A_{b,e}^c$. An interior layer $\ell\in\{1,\ldots,\lambda\}$ consists of the six vertices $(e,\ell,x,y)$ with $y\in C$ and $x\in C\setminus\{m_\ell\}$, where $m_\ell=0$ for odd $\ell$ and $m_\ell=1$ for even $\ell$. Two cells $(x,y)$ and $(x',y')$ in equal or consecutive layers of a wire are adjacent when $x\ne x'$ and $y\ne y'$; inside an end grid these are the edges of the chart. A word carries its measurement, and $(e,\ell,x,y)$ carries the measurement that returns $x$ on the plane of $e$ and $y$ on its complement.

\begin{theorem}\label{thm:FR}
Let $R\ge2$. The graph $F_R$ has $234+54(R-2)$ vertices, which carry the $165$ measurements of $F_{\rm phys}$. These measurements form a quantum $3$-coloring of $F_R$ in local dimension four whose compatibility radius is exactly $R$. There is a tower of length one over $F_R$, so $w^3\ne0$ on $\B(F_R)$ and $(\cq,\chi)=(3,5)$. No quantum $3$-coloring of $F_R$ has local dimension less than four.
\end{theorem}
\begin{proof}
The order is $6\cdot39+9\cdot6\lambda$.

\emph{Coloring and upper bound.} Sending a vertex to its measurement is a homomorphism $F_R\to F_\triangle$: the edges of a chart go to edges of $E_b$, and an edge of a wire goes to an edge of the grid of its plane, which is complete in $F_\triangle$. So the measurements form a quantum $3$-coloring and $\chi(F_R)\le5$. A chart contains the triangle of the three constant words, so $\cq(F_R)=3$.

\emph{Compatibility.} Every vertex of the chart of $b$ and every vertex of a wire at $b$, including the end grid at its other end, carries a measurement that is diagonal in basis $b$. A path between two charts runs through a wire and has at least $\lambda+1$ edges. As in the proof of~\cite[Lemma~2.2]{Z}, a path with at most $\lambda+2$ edges therefore has all its vertices in the chart of one basis $b$ or on wires at $b$, and the measurements at its ends commute.

\emph{The radius is not larger.} Take the wire of $e=24$, for which $A_{2,e}=A_2$ and $A_{4,e}=A_1$. Choose colors $z_0,\ldots,z_{\lambda+1}$ with $z_\ell\ne z_{\ell+1}$ for $0\le\ell\le\lambda$ and $z_\ell\ne m_\ell$ for $1\le\ell\le\lambda$; this is possible because every interior layer leaves two colors. The cells $(z_\ell,z_\ell)$ of the layers $\ell=0,\ldots,\lambda+1$ form a path with $\lambda+1$ edges from the constant word $z_0$ of chart $2$ to the constant word $z_{\lambda+1}$ of chart $4$. In chart $2$ the word $\omega(z_0-1,\delta_1)$ is adjacent to the constant word $z_0$, since $\delta_1$ precedes $O$, and one outcome of its measurement is the projection $Q^2_2$ onto the ray of $(0,0,1,0)$. In chart $4$ the word $\omega(z_{\lambda+1}-1,\delta_0)$ is adjacent to the constant word $z_{\lambda+1}$, and one outcome of its measurement is the projection $Q^4_1$ onto the ray of $(1,0,-1,0)$. These two projections do not commute, and the two words are joined by a path with $\lambda+3=R+1$ edges.

\emph{A tower.} We define an equivariant map $Z_{21}\to\B(F_R)$. On the copy of basis $b$ it is the map of Section~\ref{sec:165}, with the words regarded as vertices of the chart of $b$; opposite shores are completely adjacent in the chart by the argument for $F_\triangle$. It remains to extend it over the cylinders. Write $Q=K_3^{\times2}$. Two homomorphisms $f,g$ from $Q$ to a graph $G$ are \emph{compatible} if $f(q)\sim g(q')$ whenever $q\sim q'$ in $Q$. For compatible $f,g$ the maps
\[
 (\mu,\nu)\mapsto\bigl((1-a)f_*\mu+a\,g_*\mu,\ (1-a)f_*\nu+a\,g_*\nu\bigr),\qquad0\le a\le1,
\]
form an equivariant homotopy from $\B(f)$ to $\B(g)$, where $f_*$ denotes the image of a measure under $f$. Fix a wire. Two consecutive layers contain every cell of $Q$ at least once, because the end grids are complete and consecutive interior layers omit different rows. A \emph{selection} for two consecutive layers assigns to every cell of $Q$ a vertex with the same coordinates in one of the two layers. Every selection is a homomorphism $Q\to F_R$, and two selections for the same two layers are compatible. For an interior layer $\ell$ let $c^-_\ell$ and $c^+_\ell$ be the selections for the layers $\ell-1,\ell$ and for the layers $\ell,\ell+1$ that use layer $\ell$ for its six cells. They are compatible: if $q\sim q'$ and neither $c^-_\ell(q)$ nor $c^+_\ell(q')$ lies in layer $\ell$, then $q$ and $q'$ both lie in the row that layer $\ell$ omits, which is impossible for adjacent cells; otherwise the two vertices lie in equal or consecutive layers and are adjacent. Let $c_0$ and $c_{\lambda+1}$ be the selections given by the two end grids. In the sequence $c_0,c^-_1,c^+_1,c^-_2,\ldots,c^+_\lambda,c_{\lambda+1}$ consecutive members are compatible, so the homotopies above combine to an equivariant homotopy from $\B(c_0)$ to $\B(c_{\lambda+1})$.

On the torus over the loop of the end grid of $e$ in the copy of basis $i$, the map of Section~\ref{sec:165} is $\B(c_0)$ composed with a map to $\B(Q)$, namely the pair of the distributions of the cell $(\lfloor\sigma\rfloor,\lfloor\sigma+\vartheta\rfloor)$, where $\vartheta$ is the coordinate along the loop. We extend it over the cylinder by the homotopy just constructed. At the other end the result is $\B(c_{\lambda+1})$ composed with the same map to $\B(Q)$, which is the map of Section~\ref{sec:165} on the torus in the copy of basis $k$: the attaching maps of the cylinders of $Z$ are chosen in~\cite[Appendix~E]{Z} so that the words at the two ends have the same color on the plane of $e$ and the same color on its complement, also at the complemented edge. This defines an equivariant map $Z_{21}\to\B(F_R)$, and with $x_0,x_1$ a tower of length one over $F_R$. Hence $w^3\ne0$ on $\B(F_R)$ and $\chi(F_R)\ge5$ by~(T1). The assertion on the local dimension follows from~(Z5) and~(Z4) as in Theorem~\ref{thm:star}.
\end{proof}

\subsection{A physical family that starts with the 165 measurements}\label{sec:Hj}

By~\cite[Corollary~6.3]{Z} the tagged graphs $S^{j-1}_\lambda(F_{\rm phys})$ have chromatic number at least $2j+3$; their order is doubly exponential in $j$. We describe a sparse family of physical graphs with the same first level.

\paragraph{Presentations and weights.}
A measurement \emph{at level $g$} is a projective measurement with three outcomes on $(\C^4)^{\otimes g}$. For a color $c$, the \emph{padded scalar} $c$ at level $g$ is the measurement that returns $c$ with certainty; it has \emph{weight} $0$. At level one we use the $165$ measurements of $F_{\rm phys}$: the three scalar ones are the padded scalars, and the other $162$ have weight $1$. A \emph{template node} at level $g\ge2$ consists of a template at a basis $i$, with parts $D_1,\ldots,D_p$, and of a \emph{child} $N^r$ for each part, which is a padded scalar at level $g-1$, a template node at level $g-1$ or, if $g=2$, one of the $162$ other measurements. It presents the measurement
\begin{equation}\label{eq:node}
 M_a=\sum_{r=1}^p\Bigl(\sum_{s\in D_r}Q^i_s\Bigr)\otimes N^r_a .
\end{equation}
A child is \emph{internal} if it is not a padded scalar. The weight of a template node is the sum of the weights of its children, plus one if at most one child is internal. A \emph{presentation} at level $g$ is a padded scalar or a template node, or, for $g=1$, one of the $165$ measurements.

For $j\ge1$ let $H_j$ be the physical graph on all measurements at level $j$ that have a presentation of weight at most $j$. Thus $H_1=F_{\rm phys}$.

\begin{lemma}\label{lem:bottom}
Let $\mathcal B$ be a connected graph in which every vertex lies in a triangle, and let $(P_{(y,v),a})$ be a quantum $3$-coloring on $\C^d$ of the categorical product $K_3\times\mathcal B$, in which $(y,v)\sim(y',v')$ when $y\ne y'$ and $v\sim v'$. Then there is a projection $L$ that commutes with all projectors of the coloring such that on the range of $L$ the measurement at $(y,v)$ does not depend on $y$, and such that for every triangle $\mathcal T$ of $\mathcal B$ the projection $L$ is the coordinate projection of the second coordinate for the restriction of the coloring to $K_3\times\mathcal T=K_3^{\times2}$. If $L\ne0$, the measurements on its range form a quantum $3$-coloring of $\mathcal B$ in local dimension $\rank L$.
\end{lemma}
\begin{proof}
Let $v$ be a vertex and $\mathcal T$ a triangle that contains it. By~(Z2), applied to $K_3\times\mathcal T$, there is a decomposition $\C^d=V_1\oplus V_2$ that reduces the projectors at the vertices of $K_3\times\mathcal T$, such that on $V_2$ the measurement at $(y,v')$, $v'\in\mathcal T$, depends only on $v'$, and on $V_1$ only on $y$, with the edge rule of $K_3$. The operator $L_v=\sum_aP_{(0,v),a}P_{(1,v),a}$ is therefore the projection onto $V_2$. As it is defined without reference to $\mathcal T$, the space $V_2$ is the same for all triangles at $v$, and $L_v=L_{v'}$ for $v'\in\mathcal T$. Moreover $L_v$ commutes with the projectors at $(y,v)$, and on the range of $I-L_v$ the three projections $P_{(z,v),a}$, $z\in K_3$, of one color sum to the identity by~(Q2).

Let $v\sim w$ and let $\xi$ be a vector in the range of $L_v$ on which the common measurement at $(y,v)$, $y\in K_3$, returns the color $a$. For every $z$ choose $y\ne z$; then $(y,v)\sim(z,w)$ and $P_{(z,w),a}\xi=0$. So $\sum_zP_{(z,w),a}$ annihilates $\xi$. This sum commutes with $L_w$ and is the identity on the range of $I-L_w$, hence $(I-L_w)\xi=0$. Such vectors $\xi$ span the range of $L_v$, so $L_v\le L_w$, and $L_v=L_w$ by symmetry. As $\mathcal B$ is connected, all $L_v$ equal one projection $L$, which commutes with all projectors. On its range the measurement at $(y,v)$ does not depend on $y$, and for $v\sim w$ the edge rule for $(0,v)$ and $(1,w)$ gives the edge rule for $v$ and $w$.
\end{proof}

\begin{theorem}\label{thm:Hj}
Let $j\ge1$.
\begin{enumerate}
\item[(a)] The measurements that are the vertices of $H_j$ form a quantum $3$-coloring of $H_j$ in local dimension $4^j$, and $\cq(H_j)=3$.
\item[(b)] There is a tower of length $j$ over $H_j$. Hence $w^{2j+1}\ne0$ on $\B(H_j)$ and $\chi(H_j)\ge2j+3$.
\item[(c)] There is an absolute constant $c_0$ with $|V(H_j)|\le c_0^{\,j}$ and $|E(H_j)|\le c_0^{\,2j}$.
\item[(d)] Every quantum $3$-coloring of $H_j$ has local dimension divisible by $2\cdot4^{j-1}$ and at least $4^j$.
\item[(e)] For $j\ge2$, $\chi(H_j)\le2\chi(H_{j-1})-1$. The graph $H_2$ has $4{,}330{,}533$ vertices and $7\le\chi(H_2)\le9$.
\end{enumerate}
\end{theorem}
\begin{proof}
\emph{Part~\textup{(a)}.} Adjacent vertices are orthogonal at every outcome. The three padded scalars have weight $0$ and form a triangle, so~(Q3) applies.

\emph{Part~\textup{(c)}.} In a presentation, call the template nodes and the nonscalar measurements at level one \emph{internal nodes}; with the relation of being a child they form a tree. An internal node without internal children contributes one to the weight, as does an internal node with exactly one internal child, and no other node contributes. If the weight is at most $j$, the tree has at most $j$ leaves, hence at most $j-1$ nodes with two or more children, and at most $2j-1$ nodes in all. Each node is one of $162$ measurements or one of nineteen templates with at most three children, each a padded scalar or an internal node. So there are at most $c_0^{\,j}$ presentations of weight at most $j$, and the number of edges is at most the square of the number of vertices.

\emph{Part~\textup{(b)}: a tower on a space of presentations.} Let $Y_1=Z_{21}$ and $Y_g=S(Y_{g-1})$ for $g\ge2$. Let $O_1\subset Y_1$ be the orbit of the point $x_0$, which we take in the copy of basis $6$, and let $O_g\subset Y_g$ consist of the constant tuples in the block of basis $6$ on the points of $O_{g-1}$. We give weights to continuous presentations of points of $Y_g$. At level one a presentation is the point itself, of weight $0$ and not internal if it lies on $O_1$, and of weight $1$ and internal otherwise. At level $g\ge2$ a presentation is a point of $O_g$, of weight $0$ and not internal; or a basis node, with a template and children at level $g-1$, whose weight is the sum of the weights of the children plus one if at most one child is internal; or a wire node, with a first child at level $g-1$ and a background in $O_{g-1}$, whose weight is the weight of the first child plus one. Basis nodes and wire nodes are internal. Let $Y_{g,m}\subseteq Y_g$ be the set of points with a presentation of weight at most $m$. These sets are invariant under the circle action. They are compact polyhedra: regarding a point of $O_1$ also as an internal presentation of weight $1$ does not change the sets $Y_{g,m}$, because it raises the sum of the weights of the children of a node by one and lowers the added unit by at most one; and with this convention $Y_{g,m}$ is a finite union of images of products of copies of $Y_1$, of $O_1$ and of intervals under piecewise linear maps.

We claim that $Y_g$ carries a tower $x_0,\ldots,x_g$ with $x_0$ a point of $O_g$ and $x_m\in C_{2m}(Y_{g,m};\Ftwo)$ for $1\le m\le g$. For $g=1$ this is the tower on $Z_{21}$, since every point of $Y_1$ has weight at most one. For the step from $g$ to $g+1$ we take the tower of~(T2) and examine the three kinds of terms of $y_m$ in the proof of Lemma~\ref{lem:supports}. In a term of correction order zero with at least two factors of positive indices $i_\nu$, which sum to $m$, the children have weights at most $i_\nu$, and the other children are points of $O_g$. If at least two children are internal, the weight is at most $m$. If at most one is internal, one of the factors of positive index contributes a point of $O_g$, of weight $0\le i_\nu-1$, and the weight is at most $(m-1)+1$. In a term of correction order $r\ge1$ the weights of the children sum to at most $m-r$, and the weight is at most $m-r+1\le m$. The points of $BKx_{m-1}$ and $KBx_{m-1}$ are presented by wire nodes over points of $Y_{g,m-1}$, with weight at most $m$. This proves the claim. For $g=j$ it gives a tower of length $j$ on the circle space $Y_{j,j}$.

\emph{Part~\textup{(b)}: the map to $\B(H_j)$.} Let $\Phi_1:Y_1\to\B(F_{\rm phys})$ be the map of Section~\ref{sec:165} and let $\Phi_g:Y_g\to\B(S^{g-1}_0(F_{\rm phys}))$ arise from it by $g-1$ applications of~(T3) with $\lambda=0$. The $(g-1)$-fold iterate of~\eqref{eq:strategy}, starting from the measurements of $F_{\rm phys}$, assigns to every vertex of $S^{g-1}_0(F_{\rm phys})$ a measurement at level $g$, and adjacent vertices receive measurements that are orthogonal at every outcome. We show, by induction on the presentation, that if a point $z\in Y_g$ has a presentation $\pi$ of weight $m$, then the measurement of every vertex in a shore of $\Phi_g(z)$ has a presentation $\pi'$ of weight at most $m$; moreover $\pi'$ is internal if $\pi$ is, except that $\pi'$ may be a padded scalar when $\pi$ is a point of $Y_1$ off $O_1$, in which case its weight is $0=m-1$.

For points of $Y_1$ this holds because $\Phi_1$ maps $O_1$ to pairs of distributions on the scalar measurements. For a point of $O_g$, $g\ge2$, the vertices in the shores are constant words in basis $6$ on vertices in the shores one level below, by the property in~(T3), and their measurements are padded scalars. For a basis node with children $\pi_1,\ldots,\pi_p$, a vertex in a shore is a word of the chart of its basis whose entries are constant on the parts and come from the shores of the children, so by~\eqref{eq:strategy} its measurement is presented by the template node with children $\pi'_1,\ldots,\pi'_p$. If every internal child $\pi_r$ has $\pi'_r$ internal, the number of internal children is unchanged and the weight does not increase. Otherwise some weight drops by one, the added unit is at most one, and again the weight does not increase. For a wire node at the edge $e=ik$ with first child $\pi_1$, a vertex in a shore is a pair in layer $0$ or $1$ of the wire of $e$ in $S_0$, that is, a word of an end grid of $e$; by~\eqref{eq:strategy} and~\eqref{eq:planes} its measurement is $P_e\otimes N+(I-P_e)\otimes c$ with $N$ presented by $\pi'_1$ and $c$ a padded scalar. It is presented by the ternary template at basis $i$ with the part $A_{i,e}$, with child $\pi'_1$ on this part and the padded scalar $c$ on the two singletons, of weight at most the weight of $\pi_1$ plus one.

Consequently, for $z\in Y_{j,j}$ the measurements of the vertices in the two shores of $\Phi_j(z)$ are vertices of $H_j$, and those of opposite shores are adjacent in $H_j$. Taking the images of the two distributions under the map from vertices to measurements gives an equivariant map $Y_{j,j}\to\B(H_j)$. With the tower on $Y_{j,j}$ this is a tower of length $j$ over $H_j$, and~(T1) gives the rest of~(b).

\emph{Part~\textup{(d)}.} Fix a quantum $3$-coloring $(P_{M,a})$ of $H_j$ on $\C^d$. For a color $y$, a basis $i$, a position $s\in\{1,2,3\}$ and a measurement $N$ at level $g-1$, let $[i,s;y,N]$ be the measurement $Q^i_s\otimes N+(I-Q^i_s)\otimes y$ at level $g$. It is presented by a template node with a ternary template at basis $i$ that has the part $\{s\}$, with $N$ on this part and the padded scalar $y$ on the other two. For a sequence $\tau=((i_1,s_1),\ldots,(i_h,s_h))$ with $0\le h\le j-1$ put
\[
 [\tau;y,N]=[i_1,s_1;y,[i_2,s_2;y,\ldots[i_h,s_h;y,N]\ldots]]
\]
for a measurement $N$ at level $j-h$. Each node has at most one internal child, so $[\tau;y,N]$ has a presentation of weight $h$ plus the weight of $N$. In particular $F_\tau(y,x)=[\tau;y,x]$, with the padded scalar $x$, has weight $h$, and $F_\tau$ is a homomorphism $K_3^{\times2}\to H_j$. Let $\Pi_\tau$ be the coordinate projection of the coordinate $x$ for the coloring pulled back along $F_\tau$; see~(Z2). For the empty sequence $\Pi_\tau=I$.

Let $h\le j-2$. We claim that $\rank \Pi_{\tau'}=\frac14\rank \Pi_\tau$ for every sequence $\tau'=(\tau,(i,s))$. This is the step of the proof of~\cite[Theorem~4.1]{Z}, which we summarize. Inserting a ternary template node at basis $i$ with three scalar children $x_1,x_2,x_3$ as the innermost measurement gives a homomorphism $K_3^{\times4}\to H_j$ in the variables $y,x_1,x_2,x_3$, with presentations of weight $h+1$. For $x_1=x_2=x_3=x$ the measurement is $F_\tau(y,x)$, so the coordinate projections of the three parts sum to $\Pi_\tau$. The projection of a part $\{s\}$ is $\Pi_{(\tau,(i,s))}$, whichever of the two templates with this part is used, because the measurement with $x$ at position $s$ and $y$ elsewhere is the same. So the projections $T^i_s=\Pi_{(\tau,(i,s))}$, $s=1,2,3$, are pairwise orthogonal, the projection of the part $A_r$ is $D^i_r=\Pi_\tau-T^i_s-T^i_t$ for $\{r,s,t\}=\{1,2,3\}$, and with $T^i_0=\Pi_\tau-\sum_sT^i_s$ one has $D^i_r=T^i_0+T^i_r$. For an edge $e=ik$ of $K_{3,3}$ the measurements with $x$ on $A_{i,e}$ in basis $i$ and with $x$ on $A_{k,e}$ in basis $k$, and $y$ on the other positions, are equal by~\eqref{eq:planes}; they are one vertex of $H_j$. Hence $D^i_r=D^k_{r'}$ if $A_{i,e}=A_r$ and $A_{k,e}=A_{r'}$, and $D^3_3=T^6_1+T^6_2=\Pi_\tau-D^6_3$ for $e=36$. On the range of $\Pi_\tau$ the involutions $2D^i_r-\Pi_\tau$ thus satisfy the relations that~\cite{Z} derives for the tagged graphs by transport along the wires, and the argument given there with~(Z3) shows that $T^i_0,\ldots,T^i_3$ have equal rank. This proves the claim. The restriction $h\le j-2$ ensures that the inserted node lies at level at least two, where its children may be arbitrary scalars.

Applying the claim $j-1$ times gives a sequence $\tau$ of length $j-1$ with $\rank \Pi_\tau=d/4^{j-1}$; in particular $4^{j-1}$ divides $d$. For this $\tau$ and a vertex $v$ of $F_{\rm phys}$, the measurement $[\tau;y,v]$ has weight at most $j$, and $(y,v)\mapsto[\tau;y,v]$ is a homomorphism $K_3\times F_{\rm phys}\to H_j$. The graph $F_{\rm phys}$ is connected and each of its vertices lies in a triangle: the three measurements $p/0,p/1,p/2$ form a triangle; a measurement with at most two colors is adjacent to the scalar of the third color; and if the word $\omega(k,\upsilon_i)$ has the color $c\ne k$ at the position $s\ge1$, it differs at all four positions from the word with $c-1$ at $s$ and $c$ elsewhere, which is $\omega(c,\delta_{s-1})$ and has two colors. Apply Lemma~\ref{lem:bottom} to the pulled-back coloring of $K_3\times F_{\rm phys}$. For the triangle of the scalar measurements the homomorphism is $F_\tau$, so $L=\Pi_\tau$, of rank $d/4^{j-1}>0$. Hence $F_{\rm phys}$ has a quantum $3$-coloring in local dimension $d/4^{j-1}$. By~(Z7), (Z5) and~(Z4) this number is even and at least four, as in Theorem~\ref{thm:star}.

\emph{Part~\textup{(e)}.} Let $j\ge2$. A child of a template node of weight at most $j$ has weight at most $j-1$. Indeed, internal children have weight at least one; so if a child had weight at least $j$, then either a second child would be internal or the unit for a node with at most one internal child would be added, and in both cases the node would have weight at least $j+1$. So every vertex $M$ of $H_j$ is a padded scalar or has the form~\eqref{eq:node} with children that present measurements $N^r\in V(H_{j-1})$. For $N\in V(H_{j-1})$ let $D_N(M)$ be the sum of the projections $\sum_{s\in D_r}Q^i_s$ over the parts $D_r$ whose child presents the measurement $N$; for the padded scalar $c$ at level $j$ let $D_N(M)=I$ if $N$ is the padded scalar $c$ at level $j-1$ and $D_N(M)=0$ otherwise. Then $M_a=\sum_ND_N(M)\otimes N_a$. The family $(D_N(M))_N$ depends only on $M$: a real linear functional on triples of matrices that separates the finitely many elements of $V(H_{j-1})$, applied to the second tensor factor of $M$, gives a self-adjoint operator whose spectral projections are the $D_N(M)$. If $M\sim M'$ in $H_j$, then
\[
 0=\sum_a\tr(M_aM'_a)=\sum_{N,N'}\tr\bigl(D_N(M)D_{N'}(M')\bigr)\sum_a\tr(N_aN'_a),
\]
and all terms are nonnegative. So $D_N(M)D_{N'}(M')\ne0$ implies $N_aN'_a=0$ for all $a$, that is, $N\sim N'$ in $H_{j-1}$. Let $\kappa$ be a proper coloring of $H_{j-1}$ with $n=\chi(H_{j-1})$ colors, and let $\kappa(M)$ be the measurement with labels in $[n]$ that assigns the label $l$ to the sum of the $D_N(M)$ with $\kappa(N)=l$. It is a vertex of the graph $\mathcal P_n$ of~(Z1), and by the preceding observation $M\sim M'$ implies that $\kappa(M)$ and $\kappa(M')$ are orthogonal at every label. So $\kappa$ induces a homomorphism $H_j\to\mathcal P_n$, and $\chi(H_j)\le2n-1$ by~(Z1). For $j=2$ this gives $\chi(H_2)\le9$, as $\chi(H_1)=5$, and $\chi(H_2)\ge7$ by~(b).

\emph{The order of $H_2$.} Put $n=165$ and $b=162$, the numbers of all and of the nonscalar measurements at level one, and write $(n)_k=n(n-1)\cdots(n-k+1)$. A template node at level two with $k$ nonscalar children has weight $k+1$ if $k\le1$ and weight $k$ otherwise, so its weight is at most two exactly if at most two of its children are nonscalar. By the preceding paragraph, a vertex $M$ of $H_2$ is determined by the projections $D_N(M)$ of the measurements $N$ at level one that occur in it, and we count the vertices by the number of these measurements. If there is one, then $M=I\otimes N$; this is a padded scalar, or it is presented by the binary template with both children equal to a nonscalar $N$. This gives $n$ vertices. If there are two, $N$ and $N'$, then $D_N$ and $D_{N'}$ are complementary sums of parts of a template. The pairs of complementary projections that arise are the $18$ pairs $\{Q^i_s,I-Q^i_s\}$ with $s\ge1$ and the nine pairs $\{P_e,I-P_e\}$; they are $27$ distinct pairs by Section~\ref{sec:square}. In a ternary template one of the two measurements occupies two parts, so the weight is at most two exactly if $N$ or $N'$ is scalar; the binary template, which gives the pair of the plane of $e=36$, admits any two measurements. This gives $27\bigl((n)_2-(b)_2\bigr)+(b)_2$ vertices. If there are three, they occupy the three parts of one of the eighteen ternary templates, whose partitions are distinct, and at least one of them is scalar; this gives $18\bigl((n)_3-(b)_3\bigr)$ vertices. More than three cannot occur. In total
\[
 n+27\bigl((n)_2-(b)_2\bigr)+(b)_2+18\bigl((n)_3-(b)_3\bigr)=165+27\cdot978+26{,}082+18\cdot237{,}660=4{,}330{,}533 .\qedhere
\]
\end{proof}

The local dimensions of the quantum $3$-colorings of the tagged graphs $G_{j,R}$ are exactly the multiples of $4^j$~\cite[Theorem~4.1]{Z}. For $H_j$ with $j\ge2$, part~(d) leaves open whether local dimensions such as $6\cdot4^{j-1}$ occur.

\section{Facts known only by computation}\label{sec:computed}

The statements of this section were verified by programs. They have no analytical proof here, and no statement of Sections~\ref{sec:complete}--\ref{sec:small} depends on them. For each item we name the program, which lies in the folder \texttt{note/checks/computed/} of the repository unless said otherwise, and say what it checks. We write $N_3(t,R)$ for the least order of a graph with chromatic number at least $t$ that has an $R$-compatible quantum $3$-coloring, as in~\cite{Z}.

\begin{itemize}
\item[(C1)] \emph{Criticality of $F_{22}$.} Deleting any vertex or any edge of $F_{22}$ leaves a graph with a proper $3$-coloring. Program: \texttt{checks/check\_small\_graph\_22.py} of~\cite{Z}, by exhaustive search over the $22+42$ deletions.

\item[(C2)] \emph{The $678$ edges of $F_*$ are necessary.} For each edge $\epsilon$ of $F_*$ there is a coloring of the $165$ measurements with four colors for which $\epsilon$ is the only edge of $F_\triangle$ with ends of equal color. Hence every subgraph of $F_\triangle$ with chromatic number five contains all edges of $F_*$; with Theorem~\ref{thm:star}, $F_*$ is the smallest subgraph of $F_\triangle$ with chromatic number five and the only one with $678$ edges, and it becomes $4$-colorable when any edge is deleted. Moreover, there are $678$ colorings of the $165$ measurements with four colors whose sets of edges of $F_{\rm phys}$ with ends of equal color are nonempty and pairwise disjoint. Hence every subgraph of $F_{\rm phys}$ with chromatic number five has at least $678$ edges. Program: \texttt{edges\_678/check\_678.py}, which rebuilds the three graphs and checks the colorings stored in \texttt{edge\_forced\_732.txt.gz} and \texttt{edge\_packing\_843.txt.gz}.

\item[(C3)] \emph{Smaller graphs for $R=2$ and $R=3$.} There is a graph with $207$ vertices, $1{,}071$ edges and chromatic number five that has a $2$-compatible quantum $3$-coloring in local dimension four, and a graph with $279$ vertices, $1{,}449$ edges and chromatic number five that has a $3$-compatible quantum $3$-coloring in local dimension four. Each becomes $4$-colorable when any vertex is deleted. Thus $N_3(5,2)\le207$ and $N_3(5,3)\le279$, compared with the $234$ and $288$ vertices of $F_2$ and $F_3$. Both graphs consist of six charts of $39$ words of the six bases, with lists of words and rules for the edges that differ from those of Section~\ref{sec:FR}. In the first graph three cells of each end grid are shared by the two charts of a wire. In the second, three of the nine wires join their two end grids directly and through three further vertices, and the other six through a layer of six vertices as in $F_3$. Each graph uses $165$ distinct measurements, not all of which belong to $F_{\rm phys}$. The lower bound for the chromatic number follows the parity argument of~\cite[Theorem~2.4]{Z}, with one step done by exhaustive search: every proper $4$-coloring of a modified chart reads the first coordinate, in the sense of~\cite[Section~2.4]{Z}, on an odd number of its three end grids. Programs: \texttt{radius\_two\_three/verify\_square\_radius\_two\_207.py}, which constructs the first graph with its projectors and checks the parity property, commutation for all $6{,}282$ pairs of vertices at distance at most two, a proper $5$-coloring and, with the option \texttt{--check-deletions square\_radius\_two\_207.json}, the stored $4$-colorings of the $207$ vertex-deleted subgraphs; and \texttt{verify\_square\_radius\_three\_replacements.py} with \texttt{verify\_replacement\_deletions.py} for the second graph. Independently of the parity argument, a search with a SAT solver confirms that neither graph has a proper $4$-coloring (\texttt{radius\_two\_three/sat\_crosscheck.py}, which is not part of the standard run because it takes about twenty minutes).

\item[(C4)] \emph{The graph $H_2$.} The graph $H_2$ of Theorem~\ref{thm:Hj} has $2{,}116{,}320{,}480$ edges. Let $\mathcal N$ be a set of first-level measurements that is independent in $F_{\rm phys}$, and let $\mathcal D$ be a family of nonzero projections on $\C^4$ any two of which have a nonzero product. The vertices $M$ of $H_2$ for which the sum of the projections $D_N(M)$, $N\in\mathcal N$, belongs to $\mathcal D$ form an independent set of $H_2$: for two such vertices some $D_N(M)D_{N'}(M')$ with $N,N'\in\mathcal N$ is nonzero, whereas adjacency would force $N\sim N'$ by the proof of Theorem~\ref{thm:Hj}(e). For the set $\mathcal N$ of $54$ measurements and the family $\mathcal D$ stored in \texttt{level\_two/independent\_set.json}, this independent set has $1{,}407{,}626$ vertices. Deleting it leaves a graph with $2{,}922{,}907$ vertices and $732{,}610{,}036$ edges. Its chromatic number is at least six, because a proper $5$-coloring would extend to a proper $6$-coloring of $H_2$, against Theorem~\ref{thm:Hj}; and it has a quantum $3$-coloring in local dimension $16$. Programs: \texttt{level\_two/family.py}, \texttt{count\_physical\_edges.py} and \texttt{count\_remaining\_edges.py}; they count by classes of vertices, and the second compares its method with a direct enumeration on small examples. They take the graph $F_{\rm phys}$ from the file \texttt{verify\_small\_graphs.py} of an earlier supplement of the project, which is included.
\end{itemize}

\paragraph{Facts with proofs that the programs replay.}
The program \texttt{note/checks/check\_note.py} recomputes the following facts, all of which are proved above: Lemma~\ref{lem:dictator} for $n=3$ and $L\le3$, by enumeration of all proper $4$-colorings of $K_3^{\times2}$ and $K_3^{\times3}$, and its parts~(a) and~(c) for $n=4$ and $L=2$, by enumeration of the $2{,}754{,}000$ proper $6$-colorings of $K_4^{\times2}$; the lower bound of Theorem~\ref{thm:complete} for the graph $S'_0(K_3)$, by exhaustive search; the identity~\eqref{eq:two-edges} on the $22$-vertex graph; the graphs $F_{22}^{(1)}$ and $F_{22}^{(2)}$ with their orders, sizes, chromatic number and coloring; the numbers $732$ and $678$, the table of Section~\ref{sec:165} and the contractibility of the $54$ complexes $\Lambda$ in twenty orders of deletion, and $\chi(F_*)\ge5$ by exhaustive search; the graphs $F_R$ for $R=2,3,4$ with their orders, the exact compatibility radius of the coloring and a proper $5$-coloring; and the count of the vertices of $H_2$, by enumeration over smaller sets of first-level measurements. One further test was attempted and did not finish: a search for a proof that $S_0(K_4)$ has no proper $6$-coloring, which would confirm the first case of Theorem~\ref{thm:complete} beyond~\cite{Z} without Lemma~\ref{lem:bg}, was stopped after nine minutes.


\section*{Acknowledgements}
The author acknowledges support from the European Research Council (ERC Grant Agreement No.~948139) and the Deutsche Forschungsgemeinschaft (DFG, German Research Foundation) under Germany's Excellence Strategy -- EXC 2004/2 -- 390534769 (ML4Q-2). The funders had no role in the research or in the preparation of this note.

\section*{Author contributions and use of AI}
Julius A. Zeiss initiated and directed the research and takes responsibility for this note. The AI systems Claude Fable~5.2 (Anthropic), ChatGPT Astra~6.2 (OpenAI), and Gemini Neon (Google) were used to assist with mathematical exploration, the development and review of arguments, literature searches, the writing of the accompanying programs, and the preparation of the note.

\section*{Data and code availability}
The folder \texttt{note/checks/} of the repository of~\cite{Z} contains the programs named in this note. The repository is available at \url{https://github.com/JuliusAZeiss/explicit-quantum-chromatic-gap}. The program \texttt{check\_note.py} replays finite facts that are proved in Sections~\ref{sec:complete}--\ref{sec:small}; no statement of those sections depends on it. The programs and certificates in \texttt{computed/} establish the facts of Section~\ref{sec:computed}, which have no other proof.

\begingroup\small
\bibliographystyle{plainnat}
\bibliography{references}
\endgroup

\end{document}
